\subsection{外尔不等式}

在本节中，我们考虑 K = C 上的希尔伯特空间。

设 E 是希尔伯特空间。对每个 \(n \in N\) ，我们回顾希尔伯特外幂 \(\hat{\Lambda}^{n}E\) 已在 EVT, V, p. 34 中定义。对每个希尔伯特空间 F 以及每个 \(u \in \mathcal{L}(E;F)\) ，我们还定义了线性映射 \(\hat{\Lambda}^{n}u \in \mathcal{L}(\hat{\Lambda}^{n}E;\hat{\Lambda}^{n}F)\) （同前）。这些构造是函子性的：对每个希尔伯特空间 G 以及每个线性映射 \(v \in \mathcal{L}(F;G)\) ，公式

\[
\widehat {\Lambda} ^ {n} v \circ \widehat {\Lambda} ^ {n} u = \widehat {\Lambda} ^ {n} (v \circ u)
\]

成立（同前，公式 (28)）。

设 H 是 E 的闭子空间，并设 \(i_{H}\) 是从 H 到 E 中的典范单射。对 E 的每个自同态 u，我们用 \(u_{H}\) 表示 H 的自同态 \(i_{H}^{*}ui_{H}\) 。

引理 2. --- 设 \(n \in N\) 。映射 \(\hat{\Lambda}^{n}i_{H}\) 是从 \(\hat{\Lambda}^{n}H\) 到 \(\hat{\Lambda}^{n}E\) 中的线性等距映射。

设 \((e_{j})_{j\in\mathrm{J}}\) 是 H 的标准正交基，\((e_{j})_{j\in\mathrm{J}^{\prime}}\) 是 E 的标准正交基，其中 \(J\subset J^{\prime}\) 。给 \(J^{\prime}\) 赋予一个全序。元素 \(e_{j_{1}}\wedge\cdots\wedge e_{j_{n}}\) （其中 \(j_{1}<\cdots<j_{n}\) 在 \(J^{\prime}\) 中（相应地，在 J 中）变化）构成 \(\widehat{\Lambda}^{n}E\) （相应地，\(\widehat{\Lambda}^{n}H\) ）的标准正交基（EVT, V, p. 34 的命题 5）。引理得证。

在下面，我们将把 \(\hat{\Lambda}^n\mathrm{H}\) 等同于 \(\hat{\Lambda}^n\mathrm{E}\) 的一个闭子空间，借助的是映射 \(\hat{\Lambda}^n i_{\mathrm{H}}\) 。

引理 3. --- 设 F 是希尔伯特空间，且 \(u \in \mathcal{L}(\mathrm{E};\mathrm{F})\)。设 \(n \in \mathbf{N}\) 。a) 我们有 ( \(\widehat{\Lambda}^n u\) ) \(^*\) = \(\widehat{\Lambda}^n (u^*)\) ；

\begin{enumerate}
\def\labelenumi{\alph{enumi})}
\setcounter{enumi}{1}
\item
  若 F = E 且 u 是埃尔米特的（相应地，正的、正规的、酉的），则 \(\hat{\Lambda}^{n}u\) 是埃尔米特的（相应地，正的、正规的、酉的）；
\item
  我们有 \(|\widehat{\Lambda}^n u| = \widehat{\Lambda}^n |u|\) ；
\item
  设 H 是 E 的闭子空间。限制 \((\widehat{\Lambda}^{n}u)|\widehat{\Lambda}^{n}\mathrm{H}\) （即 \(\widehat{\Lambda}^{n}u\) 在 \(\widehat{\Lambda}^{n}\mathrm{H}\) 上的限制）等于 \(\widehat{\Lambda}^{n}(u|\mathrm{H})\) （等式在 \(\mathcal{L}(\widehat{\Lambda}^{n}\mathrm{H};\widehat{\Lambda}^{n}\mathrm{F})\) 中成立）；
\item
  假设 F = E。设 H 是 E 的闭子空间。我们有 \((\hat{\Lambda}^{n}u)_{\hat{\Lambda}^{n}\mathrm{H}}=\hat{\Lambda}^{n}(u_{\mathrm{H}})\) （在 \(\mathcal{L}(\hat{\Lambda}^{n}\mathrm{H})\) 中）。
\end{enumerate}

断言 \(a\)）由 EVT, V, p. 39 的公式 (13) 得出。它立即蕴含：\(\widehat{\Lambda}^{n}u\) 在 \(u\) 是埃尔米特的（相应地，酉的）时是埃尔米特的（相应地，酉的）。若 \(u\) 是正的，则 \(u^{1/2}\) 是埃尔米特的，因而 \(\widehat{\Lambda}^{n}(u^{1/2})\) 也是；关系式 \(\widehat{\Lambda}^{n}u = (\widehat{\Lambda}^{n}u^{1/2})^{2}\) 蕴含 \(\widehat{\Lambda}^{n}u\) 是正的（I, p. 138，命题 8）。这就证明了 \(b\)）。

利用前述结果，可以算得

\[
\begin{array}{r l} & {\left(\widehat {\Lambda} ^ {n} | u |\right) ^ {2} = \widehat {\Lambda} ^ {n} (| u | ^ {2}) = \widehat {\Lambda} ^ {n} (u ^ {*} u)} \\ & {\qquad = \widehat {\Lambda} ^ {n} u ^ {*} \widehat {\Lambda} ^ {n} u = \left(\widehat {\Lambda} ^ {n} u\right) ^ {*} (\widehat {\Lambda} ^ {n} u) = | \widehat {\Lambda} ^ {n} u | ^ {2}.} \end{array}
\]

由于 \(\widehat{\Lambda}^n |u|\) 依 \(b\)）是正的，故由 I, p. 118 的命题 16 得 \(\widehat{\Lambda}^n |u| = |\widehat{\Lambda}^n u|\) ，由此得 \(c\)）。

设 \(i_{\mathrm{H}}\) 是从 H 到 E 中的典范单射。空间 \(\widehat{\Lambda}^{n}\mathrm{H}\) 等同于 \(\widehat{\Lambda}^{n}\mathrm{E}\) 的一个闭子空间，借助的是映射 \(\widehat{\Lambda}^{n}i_{\mathrm{H}}\) ，由此得

\[
(\widehat {\Lambda} ^ {n} u) | \widehat {\Lambda} ^ {n} \mathrm{H} = \widehat {\Lambda} ^ {n} u \circ \widehat {\Lambda} ^ {n} i _ {\mathrm{H}} = \widehat {\Lambda} ^ {n} (u \circ i _ {\mathrm{H}}) = \widehat {\Lambda} ^ {n} (u | \mathrm{H}),
\]

以及

\[
(\widehat {\Lambda} ^ {n} u) _ {\widehat {\Lambda} ^ {n} \mathrm{H}} = (\widehat {\Lambda} ^ {n} i _ {\mathrm{H}}) ^ {*} \circ \widehat {\Lambda} ^ {n} u \circ \widehat {\Lambda} ^ {n} i _ {\mathrm{H}} = \widehat {\Lambda} ^ {n} (i _ {\mathrm{H}} ^ {*} u i _ {\mathrm{H}}) = \widehat {\Lambda} ^ {n} (u _ {\mathrm{H}}),
\]

这就证明了 \(d\)）与 \(e\)）。

设 F 是希尔伯特空间，且 \(u \in \mathcal{L}^{c}(\mathrm{E};\mathrm{F})\) 。如同 IV, p. 151 的 no. 3 中那样，我们用 \(I_{E}\) 表示 E 的满足 \(F \neq E\) 的有限维子空间 F 的维数所成的集。我们回顾，\((\alpha_{n}(u))_{n \in I_{E}}\) 表示 u 的奇异值的扩张序列（IV, p. 156 的定义 3）。

命题 10. --- 我们有等式

\[
\prod_ {i = 0} ^ {n} \alpha_ {i} (u) = \left\| \widehat {\Lambda} ^ {n + 1} u \right\|
\]

对所有 \(n \in I_{E}\) 。

引理 3 的 c) 表明 \(\|\widehat{\Lambda}^{n+1}u\|=\|\widehat{\Lambda}^{n+1}|u\|\) ；此外，由于 \(\alpha_{i}(u)=\alpha_{i}(|u|)\) 对所有 \(i\in I_{E}\) 成立（IV, p. 157 的注 5, (2)），只需对 \(|u|\) 证明断言即可。于是可设 F=E 且 u 是正的。

于是设 \(\mathrm{B}=(e_{j})_{j\in\mathrm{J}}\) 是 E 的标准正交基，使得 u 在基 B 中可对角化（IV, p. 149 的定理 1），并设 \((\lambda_{j})_{j\in\mathrm{J}}\) 是 u 在基 B 中的特征值族。给 J 赋予一个全序，并用 \(J_{n}\) 表示 J 的元素的严格递增族 \((j_{0},\ldots,j_{n})\in\mathbf{J}^{n+1}\) 所成的集。向量

\[
e _ {\iota} = e _ {j _ {0}} \wedge \dots \wedge e _ {j _ {n}}
\]

（其中 \(\iota=(j_{0},\ldots,j_{n})\in\mathrm{J}_{n}\)）构成标准正交基 \(B_{n}\) ，它是 \(\widehat{\Lambda}^{n+1}\mathrm{E}\) 的标准正交基（EVT, V, p. 34，命题 5）。对每个 \(\iota=(j_{0},\ldots,j_{n})\in\mathrm{J}_{n}\) ，记

\[
\lambda_ {\iota} = \prod_ {j = 0} ^ {n} \lambda_ {i _ {j}}.
\]

由此可得 \((\hat{\Lambda}^{n + 1}u)e_{\iota} = \lambda_{\iota}e_{\iota}\) ，因而 \(\hat{\Lambda}^{n + 1}u\) 在基 \(\mathrm{B}_n\) 中是对角的。于是，我们有

\[
\| \widehat {\Lambda} ^ {n + 1} u \| = \sup _ {\iota \in \mathrm{I} _ {n}} \lambda_ {\iota}
\]

（IV, p. 147 的引理 1），它等于 \(\lambda_{0}(u)\cdots\lambda_{n}(u)\) ，即最大的 \(n+1\) 个特征值（\(u\) 的）之积。由于 \(\lambda_{i}(u)=\alpha_{i}(u)\) 对所有 \(i\in\mathrm{I}_{\mathrm{E}}\) 在 \(u\) 为正时成立（IV, p. 157 的注 5, (3)），即得所欲证的公式。

特别地，若 \(\alpha_{1}(u) < \alpha_{0}(u) = \| u\|\) ，则可看出，不等式 \(\| \widehat{\Lambda}^{n}(u)\| \leqslant \| u\|^{n}\) （EVT, V, p. 34，公式 (29)）在 \(n \geqslant 2\) 时一般不取等号。

推论. --- 设 G 是希尔伯特空间，且 \(v \in \mathcal{L}^{\mathrm{c}}(\mathrm{F};\mathrm{G})\) 。我们有

\[
\prod_ {i = 0} ^ {n} \alpha_ {i} (v \circ u) \leqslant \prod_ {i = 0} ^ {n} \alpha_ {i} (u) \alpha_ {i} (v)
\]

对所有 \(n \in \mathrm{I}_{\mathrm{E}}\) 。

只需注意

\[
\| \widehat {\Lambda} ^ {n + 1} (v u) \| = \| \widehat {\Lambda} ^ {n + 1} v \circ \widehat {\Lambda} ^ {n + 1} u \| \leqslant \| \widehat {\Lambda} ^ {n + 1} v \| \| \widehat {\Lambda} ^ {n + 1} u \|
\]

并应用命题 10。

引理 4. --- 设 A 是环。设 \(n \in \mathbf{N}\) ，并设 \((a_i)_{0 \leqslant i \leqslant n}\) 与 \((b_i)_{0 \leqslant i \leqslant n}\) 是 A 的元素族。对 \(0 \leqslant j \leqslant n\) ，置

\[
\mathrm{A} _ {j} = \sum_ {i = 0} ^ {j} a _ {i}.
\]

我们有

\[
\sum_ {i = 0} ^ {n} a _ {i} b _ {i} = \mathrm{A} _ {n} b _ {n} - \sum_ {i = 0} ^ {n - 1} \mathrm{A} _ {i} (b _ {i + 1} - b _ {i}).
\]

置 \(A_{-1}=0\) 。我们有 \(a_{i}=A_{i}-A_{i-1}\) （\(0\leqslant i\leqslant n\)），因而

\[
\sum_ {i = 0} ^ {n} a _ {i} b _ {i} = \sum_ {i = 0} ^ {n} (\mathrm{A} _ {i} - \mathrm{A} _ {i - 1}) b _ {i} = \sum_ {i = 0} ^ {n - 1} \mathrm{A} _ {i} (b _ {i} - b _ {i + 1}) + \mathrm{A} _ {n} b _ {n},
\]

即得所欲证。

设 I 是 \(\overline{R}\) 的一个不含 \(+\infty\) 的区间。回顾（FVR, I, p. 38，注），称连续函数 \(f: I \to [-\infty, +\infty[\) 是凸的，如果它在 I 的内部上的限制是取实值的凸函数；此时它在 inf I 处有（有限或无穷的）右极限。在下面的陈述中，按约定，0 与 \(\{-\infty, +\infty\}\) 中元素之积等于 0。

引理 5. --- 设 I ⊂ R 是一个不含 +∞ 的区间，f 是从 I 到 \([-\infty, +\infty[\) 中的一个递增凸函数。

设 n 是自然数，且

\[
a _ {0} \geqslant a _ {1} \geqslant \dots \geqslant a _ {n}, \quad b _ {0} \geqslant b _ {1} \geqslant \dots \geqslant b _ {n}
\]

是 I 的元素。

设 \(\varrho_{0} \geqslant \varrho_{1} \geqslant \cdots \geqslant \varrho_{n}\) 是正实数。假设

\[
\sum_ {i = 0} ^ {j} a _ {i} \leqslant \sum_ {i = 0} ^ {j} b _ {i}.\tag{5}
\]

对所有满足 \(0 \leqslant j \leqslant n\) 的整数 j 成立。则我们有

\[
\sum_ {i = 0} ^ {n} \varrho_ {i} f (a _ {i}) \leqslant \sum_ {i = 0} ^ {n} \varrho_ {i} f (b _ {i}).\tag{6}
\]

首先假设 \(a_i \in \mathring{\mathrm{I}}\) （从而 \(f(a_i) \in \mathbf{R}\) ）对所有 \(i\) 成立。设 \(i\) 满足 \(0 \leqslant i \leqslant n\) ，并设 \((\alpha_i, \beta_i)\) 是实数，使得方程为 \(y = \alpha_i x + \beta_i\) 的直线是 \(f\) 的图形在点 \((a_i, f(a_i))\) 处的支撑直线（FVR, I, p. 37）。于是有 \(f(a_i) = \alpha_i a_i + \beta_i\) 与 \(f(b_i) \geqslant \alpha_i b_i + \beta_i\) ，因为 \(f\) 的图形位于支撑直线的上方，由此得

\[
f (b _ {i}) - f (a _ {i}) \geqslant \alpha_ {i} (b _ {i} - a _ {i}).\tag{7}
\]

此外，若 \(i < n\) ，我们有

\[
\alpha_ {i} \geqslant f _ {g} ^ {\prime} (a _ {i}) \geqslant f _ {d} ^ {\prime} (a _ {i + 1}) \geqslant \alpha_ {i + 1}
\]

（同前及 FVR, I, p. 36，推论 1）；此外 \(\alpha_{i} \geqslant 0\) ，因为 \(f\) 递增（FVR, I, p. 22，推论），从而 \(\varrho_{i}\alpha_{i} \geqslant \varrho_{i+1}\alpha_{i+1} \geqslant 0\) （\(0 \leqslant i < n\)）。

设 \(j\) 是满足 \(0 \leqslant j \leqslant n\) 的整数。置

\[
\mathrm{A} _ {j} = \sum_ {i = 0} ^ {j} (b _ {i} - a _ {i}),
\]

于是由假设 (5) 知 \(A_{j} \geqslant 0\) 。依次应用不等式 (7) 与引理 4，我们导出

\[
\begin{array}{l} \sum_ {i = 0} ^ {n} \varrho_ {i} (f (b _ {i}) - f (a _ {i})) \geqslant \sum_ {i = 0} ^ {n} \varrho_ {i} \alpha_ {i} (b _ {i} - a _ {i}) \\ \qquad = \varrho_ {n} \alpha_ {n} A _ {n} + \sum_ {j = 0} ^ {n - 1} (\varrho_ {j} \alpha_ {j} - \varrho_ {j + 1} \alpha_ {j + 1}) A _ {j} \geqslant 0. \end{array}
\]

考虑一般情形，并对 \(n\) 作归纳推理。若某个 \(a_{i}\) 不属于 I 的内部，则必有 \(a_{0} = \sup I\) 或 \(a_{n} = \inf I\) 。

先设 \(a_{n}=\inf I\) 。此时有 \(a_{n}\leqslant b_{n}\) ，从而 \(\varrho_{n}f(a_{n})\leqslant\varrho_{n}f(b_{n})\) 。把归纳假设应用于族 \((a_{0},\ldots,a_{n-1})\) 、\((b_{0},\ldots,b_{n-1})\) 与 \((\varrho_{0},\ldots,\varrho_{n-1})\) ，便得

\[
\sum_ {i = 0} ^ {n - 1} \varrho_ {i} f (a _ {i}) \leqslant \sum_ {i = 0} ^ {n - 1} \varrho_ {i} f (b _ {i}),
\]

再加上 \(\varrho_{n}f(a_{n})\)即得所欲证的不等式。\(a_{0} = \sup I\) 的情形可类似处理。

命题 11（外尔不等式）. --- 设 G 是希尔伯特空间，且 \(v \in \mathcal{L}^{\mathrm{c}}(\mathrm{F};\mathrm{G})\)。设 \(g: \mathbf{R}_{+} \to [-\infty, +\infty[\) 是不减函数，使得函数 \(g \circ \exp\) 是凸的。我们有

\[
\sum_ {i = 0} ^ {n} g \left(\alpha_ {i} (v \circ u)\right) \leqslant \sum_ {i = 0} ^ {n} g \left(\alpha_ {i} (v) \alpha_ {i} (u)\right)
\]

对所有 \(n \in I_{E} \cap I_{F}\) 。

置 \(I = [-\infty, +\infty[\) 且 \(f = g \circ \exp\)。可以对下列取值应用引理 5

\[
a _ {i} = \log (\alpha_ {i} (v \circ u)) \in \mathrm{I}, \quad b _ {i} = \log (\alpha_ {i} (v) \alpha_ {i} (u)) \in \mathrm{I}
\]

并取 \(\varrho_{i} = 1\) （\(0\leqslant i\leqslant n\)），因为

\[
\sum_ {i = 0} ^ {j} a _ {i} = \log \Bigl (\prod_ {i = 0} ^ {j} \alpha_ {i} (v \circ u) \Bigr) \leqslant \log \Bigl (\prod_ {i = 0} ^ {j} \alpha_ {i} (v) \alpha_ {i} (u) \Bigr) = \sum_ {i = 0} ^ {j} b _ {i}
\]

对所有 \(0 \leqslant j \leqslant n\) 成立，依据是命题 10 的推论。于是不等式 (6) 就是所要的结论。

引理 6. --- 设 g 与 h 分别是定义在 R 的区间 I 与 J 上的凸函数。若 g 递增且定义在 h 的像上，则函数 g ∘ h 在 J 上是凸的。

事实上，对 \(t \in [0,1]\) 与 \((x,y) \in \mathrm{J} \times \mathrm{J}\) ，我们有

\[
\begin{array}{c} g (h (t x + (1 - t) y)) \leqslant g (t h (x) + (1 - t) h (y)) \\ \leqslant t g (h (x)) + (1 - t) g (h (y)). \end{array}
\]

推论. --- 设 G 是希尔伯特空间，且 \(v \in \mathcal{L}^{\mathrm{c}}(\mathrm{F};\mathrm{G})\)。设 \(n \in I_{E} \cap I_{F}\) 。

\begin{enumerate}
\def\labelenumi{\alph{enumi})}
\tightlist
\item
  设 \(r \in \mathbf{R}_{+}^{*}\) 。我们有
\end{enumerate}

\[
\sum_ {i = 0} ^ {n} \alpha_ {i} (v \circ u) ^ {r} \leqslant \sum_ {i = 0} ^ {n} \alpha_ {i} (v) ^ {r} \alpha_ {i} (u) ^ {r}.
\]

\begin{enumerate}
\def\labelenumi{\alph{enumi})}
\setcounter{enumi}{1}
\tightlist
\item
  设 \(p, q, r \in \mathbf{R}_{+}^{*}\) 满足 \(\frac{1}{p} + \frac{1}{q} = \frac{1}{r}\) 。则
\end{enumerate}

\[
\left(\sum_ {i = 0} ^ {n} \alpha_ {i} (v \circ u) ^ {r}\right) ^ {1 / r} \leqslant \left(\sum_ {i = 0} ^ {n} \alpha_ {i} (v) ^ {p}\right) ^ {1 / p} \left(\sum_ {i = 0} ^ {n} \alpha_ {i} (u) ^ {q}\right) ^ {1 / q}.
\]

\begin{enumerate}
\def\labelenumi{\alph{enumi})}
\setcounter{enumi}{2}
\tightlist
\item
  假设 F = E。对每个整数 \(m \geqslant 2\) ，我们有
\end{enumerate}

\[
\sum_ {i = 0} ^ {n} \alpha_ {i} (u ^ {m}) \leqslant \left(\sum_ {i = 0} ^ {n} \alpha_ {i} (u) ^ {2}\right) ^ {m / 2}.
\]

设 \(r \in R_{+}^{*}\) ，并设 g 是在 \(\mathbf{R}_{+}\) 上由 \(g(x) = x^{r}\) 定义的函数。函数 \(g \circ \exp\) 是凸的（引理 6），故把命题 11 应用于函数 g 即得断言 a)。

断言 \(b\)）由 \(a\)）及赫尔德不等式（INT, I，命题 4）得出。

设 \(m \geqslant 2\) 是整数。我们把 \(b)\) 应用于 \(r = 1\) 、\(p = q = 2\) 与 \(v = u^{m-1}\) ，得到

\[
\sum_ {i = 0} ^ {n} \alpha_ {i} (u ^ {m}) \leqslant \left(\sum_ {i = 0} ^ {n} \alpha_ {i} (u ^ {m - 1}) ^ {2}\right) ^ {1 / 2} \left(\sum_ {i = 0} ^ {n} \alpha_ {i} (u) ^ {2}\right) ^ {1 / 2}.
\]

我们对 \(m \geqslant 2\) 作归纳来证明 c)。前一不等式给出了 m = 2 时的断言。设 \(m \geqslant 3\) 且断言对 m - 1 成立；由于我们有

\[
\sum_ {i = 0} ^ {n} \alpha_ {i} (u ^ {m - 1}) ^ {2} \leqslant \left(\sum_ {i = 0} ^ {m} \alpha_ {i} (u ^ {m - 1})\right) ^ {2},
\]

上面的不等式蕴含

\[
\sum_ {i = 0} ^ {n} \alpha_ {i} (u ^ {m}) \leqslant \left(\sum_ {i = 0} ^ {n} \alpha_ {i} (u ^ {m - 1})\right) ^ {1 / 2} \left(\sum_ {i = 0} ^ {n} \alpha_ {i} (u) ^ {2}\right) ^ {1 / 2} \leqslant \left(\sum_ {i = 0} ^ {n} \alpha_ {i} (u) ^ {2}\right) ^ {m / 2}
\]

其中应用了归纳假设。

\subsection{有限迹自同态}

我们回顾，E 的正自同态 \(u\) 是有限迹的，当且仅当存在 E 的标准正交基 \((e_i)_{i \in I}\) ，使得

\[
\sum_ {i \in \mathrm{I}} \langle e _ {i} | u (e _ {i}) \rangle <   + \infty
\]

（EVT, V, p. 48，引理 3 及 p. 49，定义 7）。若 K = C，E 的有限迹自同态所成的空间 \(\mathcal{L}_1(\mathrm{E})\) 是由正的有限迹自同态组成的集生成的向量空间（EVT, V, p. 50，定义 8）；若 K = R，空间 \(\mathcal{L}_1(\mathrm{E})\) 定义为交 \(\mathcal{L}(\mathrm{E}) \cap \mathcal{L}_1(\mathrm{E}_{(\mathbf{C})})\) （EVT, V, p. 50）。

若 \(u\in \mathcal{L}_1(\mathrm{E})\) ，则级数

\[
\sum_ {i \in \mathrm{I}} \langle e _ {i} | u (e _ {i}) \rangle
\]

对 E 的每个标准正交基 \((e_{i})_{i\in\mathrm{I}}\) 收敛，且其和与标准正交基无关；这个和称为 u 的迹 \(\operatorname{Tr}(u)\) （EVT, V, p. 50）。若 K = R，我们有 \(\operatorname{Tr}(u) = \operatorname{Tr}(u_{(\mathbf{C})})\) 。

设 \(u \in \mathcal{L}_1(\mathrm{E})\) 。我们有 \(u^* \in \mathcal{L}_1(\mathrm{E})\) ，且 \(\operatorname{Tr}(u^*) = \overline{\operatorname{Tr}(u)}\) （同前）。

命题 12. --- 设 \(B = (e_i)_{i \in I}\) 是 E 的标准正交基。设 \(u \in \mathcal{D}_B(E)\) ，并用 \(\lambda = (\lambda_i)_{i \in I}\) 表示它的特征值族。自同态 \(u\) 是有限迹的，当且仅当族 \(\lambda\) 在 K 中可和。此时我们有 \(\operatorname{Tr}(u) = \sum \lambda_i\) 。

如有必要，把 \(u\) 换成 \(u_{(\mathbf{C})}\) ，我们可设 \(K = \mathbf{C}\) 。

先设 \(u\) 是有限迹的。根据 EVT, V, p. 50 及 p. 49 的公式 (25)，族 \((\langle e_i | u(e_i) \rangle)_{i \in I} = (\lambda_i)_{i \in I}\) 是可和的。

反之，设族 \(\lambda\) 是可和的。

于是族 \((\mathcal{R}(\lambda_{i})^{+})\) 、\((\mathcal{R}(\lambda_{i})^{-})\) 、\((\mathcal{I}(\lambda_{i})^{+})\) 、\((\mathcal{I}(\lambda_{i})^{-})\) 都是可和的。自同态 u 是 \(\mathcal{D}_{\mathrm{B}}(\mathrm{E})\) 中以这些族为特征值的元素的线性组合。\(\mathcal{D}_{\mathrm{B}}(\mathrm{E})\) 的这些元素都是正的。由于按定义，有限迹自同态空间由正的有限迹自同态生成，因此可设对所有 i 有 \(\lambda_{i} \geqslant 0\) 。由于 \(\langle e_{i} | u(e_{i}) \rangle = \lambda_{i}\)，族 \((\langle e_{i} | u(e_{i}) \rangle)_{i \in \mathrm{I}}\) 是可和的，故 u 是有限迹的（EVT, V, p. 48，引理 3）。

最后，若 u 有有限迹，则由 EVT, V, p. 50，我们有

\[
\operatorname{Tr} (u) = \sum_ {i \in \mathrm{I}} \left\langle e _ {i} \mid u \left(e _ {i}\right) \right\rangle = \sum_ {i \in \mathrm{I}} \lambda_ {i}.
\]

推论 1. --- 设 u 是 E 的有限迹自同态。

\begin{enumerate}
\def\labelenumi{\alph{enumi})}
\item
  自同态 u 是紧的；
\item
  设 \(\mathrm{B} = (e_i)_{i \in \mathrm{I}}\) 是 u 的始空间 \(\operatorname{Ker}(u)^{\circ}\) 的标准正交基，\(\mathrm{C} = (f_i)_{i \in \mathrm{I}}\) 是 E 的标准正交族，且 \((\alpha_i)_{i \in \mathrm{I}}\) 是 \((\mathbf{R}_+^*)^{\mathrm{I}}\) 中的族，使得 \(u(e_i) = \alpha_i f_i\) 对所有 \(i \in I\) 成立（IV, p. 150 的推论 2）。我们有
\end{enumerate}

\[
\operatorname{Tr} (u) = \sum_ {i \in I} \alpha_ {i} \langle e _ {i} | f _ {i} \rangle .
\]

为证明 \(a\)），可设 \(\mathbf{K} = \mathbf{C}\) （EVT, V, p. 50 及 III, p. 2 的注 4），且 \(u\) 是正的（EVT, V, p. 50，定义 8）。

由 EVT, V, p. 56 的推论 1，存在 E 的标准正交基 \(B = (e_i)_{i \in I}\) ，使得 \(u\) 在基 B 中是对角的，且特征值族 \(\lambda\) （\(u\) 的特征值族）属于 \(\mathbf{R}_+^I\) 并可和。因此族 \(\lambda\) 属于 \(\mathcal{C}_0(\mathrm{I};\mathrm{K})\) （TG, III, p. 38，命题 1），从而自同态 \(u\) 是紧的（IV, p. 148 的命题 2）。

我们来证明 \(b\)）。设 \((e_i)_{i \in \mathbf{J}}\) 是 E 的标准正交基，它扩张族 B，使得 \(u(e_i) = 0\) （当 \(i \in \mathbf{J} - \mathbf{I}\) 时）。我们得到

\[
\operatorname{Tr} (u) = \sum_ {i \in \mathrm{J}} \langle e _ {i} \mid u (e _ {i}) \rangle = \sum_ {i \in \mathrm{I}} \langle e _ {i} \mid u (e _ {i}) \rangle = \sum_ {i \in \mathrm{I}} \alpha_ {i} \langle e _ {i} \mid f _ {i} \rangle .
\]

推论 2. --- 设 F 是希尔伯特空间。设 \(u \in \mathcal{L}(\mathrm{E};\mathrm{F})\) 是希尔伯特--施密特算子。则 u 是紧线性算子。

设 \((j,|u|)\) 是 u 的极分解（I, p. 140，定义 4）。按定义（EVT, V, p. 50，定义 9），自同态 \(u^{*}u\) 有有限迹，因而是紧的（推论 1, a)）；对 \(|u|=\sqrt{u^{*}u}\) （III, p. 91 的命题 6, b)）与 \(u=j|u|\) （III, p. 5 的命题 3）结论同样成立。

推论 3. --- 假设 K = C。设 u 是 E 的正紧自同态。自同态 u 是有限迹的，当且仅当它的特征值的递减族 \((\lambda_{n}(u))_{n\in\mathrm{I}_{\mathrm{E}}}\) 可和；此时 u 的迹就是这个族的和。

由 IV, p. 149 的定理 1，这可由前一命题以及序列 \((\lambda_n(u))_{n\in \mathrm{I_E}}\) 的定义（IV, p. 151 的 no. 3）得出，并用到 EVT, V, p. 49 的公式 (28)。

推论 4. --- 设 F 是希尔伯特空间。设 \(u \in \mathcal{L}(\mathrm{E};\mathrm{F})\) 是紧线性映射。设 \(\mathrm{B} = (e_i)_{i \in \mathrm{I}}\) 是始空间 \(\operatorname{Ker}(u)^{\circ}\) （\(u\) 的始空间）的标准正交基，\(\mathrm{C} = (f_i)_{i \in \mathrm{I}}\) 是 F 的标准正交族，且 \((\alpha_{i})_{i\in\mathrm{I}}\) 是 \((\mathbf{R}_{+}^{*})^{\mathrm{I}}\) 中的族，使得 \(u(e_{i})=\alpha_{i}f_{i}\) 对所有 \(i\in I\) 成立（IV, p. 150 的推论 2）。

E 的自同态 \(|u|\) 有有限迹，当且仅当族 \((\alpha_{i})_{i\in \mathrm{I}}\) 可和。此时我们有 \(u\in \mathcal{L}_2(\mathrm{E};\mathrm{F})\) ，并且

\[
\operatorname{Tr} (| u |) = \sum_ {i \in \mathrm{I}} \alpha_ {i}, \quad \| u \| _ {2} ^ {2} = \operatorname{Tr} (u ^ {*} u) = \sum_ {i \in \mathrm{I}} \alpha_ {i} ^ {2},\tag{8}
\]

特别地有 \(\|u\|_{2}\leqslant\mathrm{Tr}(|u|)\) 。

\(|u|\) 的非零特征值族是 \((\alpha_{i})_{i\in \mathrm{I}}\) ，故 \(|u|\) 有有限迹当且仅当族 \((\alpha_{i})_{i\in \mathrm{I}}\) 可和（命题 12）。若确实如此，非零特征值族 \((\alpha_{i}^{2})\) （\(u^{*}u = |u|^{2}\) 的非零特征值族）可和，且公式 (8) 由同前得出。

引理 7. --- 设 \(\lambda = (\lambda_i)_{i \in \mathrm{I}}\) 是复数族。对每个 \(t \in \mathbf{C}^*\) ，用 \(n_t\) 表示 \(i \in \mathrm{I}\) 中使 \(\lambda_i = t\) 者所成之集的基数。族 \((\lambda_i)_{i \in \mathrm{I}}\) 可和，当且仅当 \(n_t\) 对所有 \(t \in \mathbf{C}^*\) 有限，且族 \((n_t t)_{t \in \mathbf{C}^*}\) 可和。此时这两个族的和相等。

设族 \((\lambda_{i})_{i\in\mathrm{I}}\) 可和。对每个 \(t\in\mathbf{C}\) ，用 \(I_{t}\) 表示 \(i\in I\) 中使 \(\lambda_{i}=t\) 者所成之集。把 TG, III, p. 39 的定理 2 应用于 I 被诸集 \(I_{t}\) 所作的分划，可知集 \(I_{t}\) 对每个 \(t\in\mathbf{C}^{*}\) 有限，且族 \((n_{t}t)_{t\in\mathbf{C}^{*}}\) 可和，其和等于族 \((\lambda_{i})_{i\in\mathrm{I}}\) 的和。

反之，设 \(n_{t}\) 对所有 \(t \in \mathbf{C}^{*}\) 有限，且族 \((n_{t}t)_{t \in \mathbf{C}^{*}}\) 可和。设 J 是 I 的有限子集，\(\Lambda\) 是诸 \(\lambda_{i}\) （\(i \in J\)）所成之集。我们有

\[
\left| \sum_ {i \in J} \lambda_ {i} \right| \leqslant \sum_ {t \in \Lambda - \{0 \}} n _ {t} | t | \leqslant \sum_ {t \in \mathbf {C} ^ {*}} n _ {t} | t |,
\]

因此族 \((\lambda_{i})_{i\in\mathrm{I}}\) 可和（TG, VII, p. 17，推论）。

命题 13. --- 假设 K = C。设 u 是 E 的紧正规自同态。对 \(t \in \mathrm{Sp}_{\mathrm{s}}(u)\) ，用 \(n_{t} \geqslant 1\) 表示 u 的特征值 t 的谱重数。u 有有限迹的必要充分条件是族 \((n_{t}t)_{t \in \mathrm{Sp}_{\mathrm{s}}(u)}\) 可和。此时 u 的迹就是这个族的和。

设 \(\mathrm{B} = (e_i)_{i \in \mathrm{I}}\) 是 E 的标准正交基，使得 u 在基 B 中是对角的（IV, p. 149 的定理 1），并设 \(\lambda = (\lambda_i)_{i \in \mathrm{I}}\) 是它的特征值族。由于对 \(t \in \mathrm{Sp}_\mathrm{s}(u)\) ，重数 \(n_t\) 等于元素 \(i \in I\) 中使 \(\lambda_i = t\) 者的个数，且 \(\lambda\) 的非零元素属于 \(\mathrm{Sp}_{\mathrm{s}}(u)\) ，断言即由命题 12 与前一引理得出。

\subsection{核型映射}

在本节中，F 表示 K 上的希尔伯特空间。当 K = C 时，我们回顾（I, p. 139，定义 3），对 \(u \in \mathcal{L}(\mathrm{E}; \mathrm{F})\) ，用 \(|u|\) 表示 E 的正自同态 \(\sqrt{u^{*} \circ u}\) 。当 K = R 时，元素 \(|u_{(\mathbf{C})}|\) （\(\mathcal{L}(\mathrm{E}_{(\mathbf{C})})\) 中的元素）具有形式 \(v_{(\mathbf{C})}\) ，其中自同态 \(v \in \mathcal{L}(\mathrm{E})\) 是唯一的且仍记作 \(|u|\) （参见 I, p. 87）。

我们用 \((u,v)\mapsto\langle u|v\rangle=\mathrm{Tr}(u^{*}v)\) 表示希尔伯特空间 \(\mathcal{L}_{2}(\mathrm{E};\mathrm{F})\) 中的内积（EVT, V, p. 53，注 1 与注 2）。

对每个 \(u \in \mathcal{L}(\mathrm{E};\mathrm{F})\) ，用 \(\| u\| _1\) 表示量 \(\mathrm{Tr}(|u|)\) （若 \(|u|\) 有有限迹）；否则置 \(\| u\| _1 = +\infty\) 。于是我们有 \(\| u\| _1 \in \mathbf{R}_+ \cup \{+\infty\}\) 。由于 \(\| u\| = \| |u|\|\) （I, p. 139 的命题 10, a)），且当 \(v\) 是正的并有限迹时有 \(\mathrm{Tr}(v) \geqslant \| v\|\) （EVT, V, p. 49，公式 (24bis) 及 p. 44，命题 9），由此可得

\[
\| u \| \leqslant \| u \| _ {1}.\tag{9}
\]

若 \(\| u\| _1\) 有限，则 \(u\) 是希尔伯特--施密特算子，且 \(\| u\| _2 \leqslant \| u\| _1\) （IV, p. 165 的推论 4）。

命题 14. --- 设 \(u \in \mathcal{L}_2(\mathrm{E};\mathrm{F})\) 。我们有

\[
\| u\|_{1} = \sup_{\substack{v\in \mathcal{L}_{2}(\mathrm{E};\mathrm{F})\\ \| v\| \leqslant 1}}|\langle v  |  u\rangle |,
\]

其中的上确界在 \(\mathbf{R}_{+} \cup \{+\infty\}\) 中计算。

设 \(\mathrm{B} = (e_i)_{i \in \mathrm{I}}\) 是 u 的始空间 \(\operatorname{Ker}(u)^{\circ}\) 的标准正交基，\(\mathrm{C} = (f_i)_{i \in \mathrm{I}}\) 是 F 中的标准正交族，\((\alpha_i)_{i \in \mathrm{I}}\) 是 \((\mathbf{R}_+^*)^{\mathrm{I}}\) 中的族，使得 \(u(e_i) = \alpha_i f_i\) 对所有 \(i \in I\) 成立（IV, p. 150 的推论 2）。族 \((\alpha_i)_{i \in \mathrm{I}}\) 平方可积，因为 u 属于 \(\mathcal{L}_2(\mathrm{E};\mathrm{F})\) （IV, p. 165 的推论 4）。设 \((e_j)_{j \in \mathrm{J}}\) 是 \((e_i)_{i \in \mathrm{I}}\) 的扩张，即 E 的标准正交基。

设 L 是 I 的有限子集。设 \(v_{L}\) 是由

\[
v _ {\mathrm{L}} (x) = \sum_ {i \in \mathrm{L}} \langle e _ {i} | x \rangle f _ {i}
\]

对所有 \(x\) 属于 E 定义的从 E 到 F 中的连续有限秩线性映射。我们有 \(\| v_{\mathrm{L}}\| \leqslant 1\) 且 \(v_{\mathrm{L}}\in \mathcal{L}_{2}(\mathrm{E};\mathrm{F})\) 。此外，对每个 \(j\in \mathrm{J}\) ，有 \(v_{\mathrm{L}}(e_j) = 0\) （若 \(j\notin \mathrm{L}\) ），且有 \(v_{\mathrm{L}}(e_j) = f_j\) （若 \(j\in \mathrm{L}\) ）。于是

\[
| \langle v _ {\mathrm{L}} | u \rangle | = | \operatorname{Tr} (v _ {\mathrm{L}} ^ {*} u) | = \left| \sum_ {j \in \mathrm{J}} \langle v _ {\mathrm{L}} (e _ {j}) | u (e _ {j}) \rangle \right| = \sum_ {j \in \mathrm{L}} \alpha_ {j}.
\]

由同前，我们推出

\[
\| u\|_{1} = \sup_{\text{L}}\sum_{j\in \text{L}}\alpha_{j}\leqslant \sup_{\substack{v\in \mathcal{L}_{2}(\text{E};\text{F})\\ \| v\| \leqslant 1}}|\langle v\mid u\rangle |\tag{10}
\]

在 \(\mathbf{R}_{+}\cup\{+\infty\}\) 中成立。

当 \(\|u\|_{1}=+\infty\) 时，这即给出命题中的等式。

设 \(\|u\|_{1}\) 有限。对每个 \(i \in I\) ，设 \(p_{i}\) 是从 E 到 F 的线性映射，满足 \(p_{i}(e_{i}) = f_{i}\) 且 \(p_{i}(x) = 0\) 对所有 \(x \in e_{i}^{\circ}\) 成立。对 I 中所有的 j 和 k，我们有

\[
\langle p _ {j} \mid p _ {k} \rangle = \mathrm{Tr} (p _ {j} ^ {*} p _ {k}) = \sum_ {i \in \mathrm{J}} \langle p _ {j} (e _ {i}) \mid p _ {k} (e _ {i}) \rangle ,
\]

除非 \(j = k\) ，它为零；在此情形这个量等于 \(\| f_k\|^2 = 1\) 。因此族 \((p_i)_{i \in \mathrm{I}}\) 在 \(\mathcal{L}_2(\mathrm{E};\mathrm{F})\) 中标准正交。从而族 \((\alpha_i p_i)_{i \in \mathrm{I}}\) 在 \(\mathcal{L}_2(\mathrm{E};\mathrm{F})\) 中可和；其和等于 \(u\) ，因为这两个连续线性映射在元素 \(e_i\) 上（对所有 \(i \in \mathrm{J}\)）相合。

设 \(v \in \mathcal{L}_2(\mathrm{E};\mathrm{F})\) 。对每个 \(i \in \mathrm{I}\) ，由此可得

\[
\langle v \mid p _ {i} \rangle = \operatorname{Tr} (v ^ {*} p _ {i}) = \sum_ {j \in \mathrm{J}} \langle v (e _ {j}) \mid p _ {i} (e _ {j}) \rangle = \langle v (e _ {i}) \mid f _ {i} \rangle .
\]

若 \(\|v\| \leqslant 1\) ，我们有

\[
\begin{array}{c} | \langle v | u \rangle | = \left| \langle v | \sum_ {i \in I} \alpha_ {i} p _ {i} \rangle \right| \leqslant \sum_ {i \in I} \alpha_ {i} | \langle v | p _ {i} \rangle | \\ = \sum_ {i \in I} \alpha_ {i} | \langle v (e _ {i}) | f _ {i} \rangle | \leqslant \sum_ {i \in I} \alpha_ {i} = \mathrm{Tr} (| u |), \end{array}
\]

由此得到不等式

\[
\sup_{\substack{v\in \mathcal{L}_{2}(\mathrm{E};\mathrm{F})\\ \| v\| \leqslant 1}}|\langle v\mid u\rangle |\leqslant \| u\|_{1},
\]

它与公式 (10) 相结合，便完成了命题的证明。

由这个命题可得，集 \(\mathcal{L}_1(\mathrm{E};\mathrm{F})\) （从 E 到 F 的连续线性映射 \(u\) 中使 \(\| u\| _1\) 有限者所成之集）是 \(\mathcal{L}_{2}(\mathrm{E};\mathrm{F})\) 的向量子空间，且映射 \(u\mapsto \| u\| _1\) 是 \(\mathcal{L}_1(\mathrm{E};\mathrm{F})\) 上的半范数；不等式 (9) 表明它是一个范数。

定义 4. --- 我们称向量空间 \(\mathcal{L}_{1}(\mathrm{E};\mathrm{F})\) （赋以范数 \(u\mapsto \| u\|_{1}\)）为从 E 到 F 的核型映射所成的空间。若 \(u\in \mathcal{L}_1(\mathrm{E};\mathrm{F})\) ，则称 \(u\) 是核型的。

注记. --- 1) 假设 K = R。设 \(u \in \mathcal{L}(E; F)\) 。我们有 \(u \in \mathcal{L}_{1}(E; F)\) 当且仅当 \(u_{(\mathbf{C})} \in \mathcal{L}_{1}(E_{(\mathbf{C})}; F_{(\mathbf{C})})\) ；此时有 \(\|u\|_{1} = \|u_{(\mathbf{C})}\|_{1}\) 。

\begin{enumerate}
\def\labelenumi{\arabic{enumi})}
\setcounter{enumi}{1}
\tightlist
\item
  设 \(u\) 是从 E 到 F 的核型映射。由于映射 \(u\) 是希尔伯特--施密特的，它是紧的（IV, p. 165 的推论 2）。此外，命题 14 蕴含对每个 \(v \in \mathcal{L}_2(\mathrm{E};\mathrm{F})\) ，我们有
\end{enumerate}

\[
| \langle v | u \rangle | \leqslant \| v \| \| u \| _ {1}.\tag{11}
\]

\begin{enumerate}
\def\labelenumi{\arabic{enumi})}
\setcounter{enumi}{2}
\item
  假设 \(\mathrm{K} = \mathbf{C}\)。设 \(u \in \mathcal{L}_1(\mathrm{E};\mathrm{E})\) 且 \(u\) 是正的。范数 \(\| u\| _1\) （\(u\) 的范数）是序列 \((\lambda_n(u))_{n\in \mathrm{I_E}}\) 的和（IV, p. 165 的推论 3）。
\item
  \(\mathcal{L}_1(\mathrm{E};\mathrm{F})\) 到 \(\mathcal{L}(\mathrm{E};\mathrm{F})\) 中的典范单射的范数 \(\leqslant 1\) （p. 167 的不等式 (9)）。
\end{enumerate}

命题 15. --- 设 G 是 K 上的希尔伯特空间。映射 \((u,v)\mapsto v\circ u\) （从 \(\mathcal{L}(\mathrm{E};\mathrm{F})\times\mathcal{L}(\mathrm{F};\mathrm{G})\) 到 \(\mathcal{L}(\mathrm{E};\mathrm{G})\) 中的映射）过渡到商，定义了范数 \(\leqslant1\) 的从 \(\mathcal{L}_{2}(\mathrm{E};\mathrm{F})\times\mathcal{L}_{2}(\mathrm{F};\mathrm{G})\) 到 \(\mathcal{L}_{1}(\mathrm{E};\mathrm{G})\) 中的连续双线性映射。

设 \(u \in \mathcal{L}_2(\mathrm{E};\mathrm{F})\) 且 \(v \in \mathcal{L}_2(\mathrm{F};\mathrm{G})\)。设 \(w \in \mathcal{L}_2(\mathrm{E};\mathrm{G})\) 。我们得到 \(\langle uv | w \rangle = \operatorname{Tr}(v^* u^* w) = \langle v | u^* w \rangle\) ，由此

\[
| \langle u v | w \rangle | \leqslant \| v \| _ {2} \| u ^ {*} w \| _ {2} \leqslant \| v \| _ {2} \| u \| _ {2} \| w \|
\]

此处用到柯西--施瓦茨不等式及 EVT, V, p. 52 的公式 (37)。因此结论由命题 14 得出。

引理 8. --- 映射 \(u \mapsto u^{*}\) （从 \(\mathcal{L}(\mathrm{E};\mathrm{F})\) 到 \(\mathcal{L}(\mathrm{F};\mathrm{E})\) 中的映射）过渡到商，定义了从 \(\mathcal{L}_{1}(\mathrm{E};\mathrm{F})\) 到 \(\mathcal{L}_{1}(\mathrm{F};\mathrm{E})\) 中的等距线性映射。

设 \(u \in \mathcal{L}_1(\mathrm{E};\mathrm{F})\) 。由于 \(u\) 是希尔伯特--施密特映射，\(u^*\) 也是（EVT, V, p. 54）。映射 \(v \mapsto v^*\) 是从 \(v \in \mathcal{L}_2(\mathrm{E};\mathrm{F})\) 中满足 \(\|v\| \leqslant 1\) 者所成之集到 \(w \in \mathcal{L}_2(\mathrm{F};\mathrm{E})\) 中满足 \(\|w\| \leqslant 1\) 者所成之集上的双射；对每个 \(v \in \mathcal{L}_2(\mathrm{E};\mathrm{F})\)

满足 \(\|v\|\leqslant1\) 者，我们有 \(\langle v|u\rangle=\langle u^{*}|v^{*}\rangle\) （EVT, V, p. 54，公式 (42)），于是结论由命题 14 得出。

命题 16. --- 设 \(E_{1}\) 与 \(F_{1}\) 是希尔伯特空间。设 \(u\) 属于 \(\mathcal{L}_{1}(\mathrm{E};\mathrm{F})\) ，\(v\) 属于 \(\mathcal{L}(\mathrm{E}_{1};\mathrm{E})\) ，\(v_{2}\) 属于 \(\mathcal{L}(\mathrm{F};\mathrm{F}_{1})\)。则有 \(v_{2}uv_{1} \in \mathcal{L}_{1}(\mathrm{E}_{1};\mathrm{F}_{1})\) 且 \(\|v_{2}uv_{1}\|_{1} \leqslant \|v_{2}\|\|v_{1}\|\|u\|_{1}\) 。

设 \(w \in \mathcal{L}_2(\mathrm{E};\mathrm{F}_1)\) 满足 \(\| w\| \leqslant 1\) 。我们有 \(v_{2}u \in \mathcal{L}_{2}(\mathrm{E};\mathrm{F}_{1})\) （EVT, V, p. 52，公式 (36)）。由于 \(v_{2}^{*}w \in \mathcal{L}_{2}(\mathrm{E};\mathrm{F})\) （同前），由此可得

\[
| \langle w | v _ {2} u \rangle | = | \langle v _ {2} ^ {*} w | u \rangle | \leqslant \| v _ {2} \| \| u \| _ {1}
\]

（公式 (11)），于是 \(v_2u \in \mathcal{L}_1(\mathrm{E};\mathrm{F}_1)\) 且 \(\| v_2u\| _1 \leqslant \| v_2\| \| u\| _1\) （命题 14）。设 \(v_{1} \in \mathcal{L}(\mathrm{E}_{1};\mathrm{E})\) ；由于 \(uv_{1} = (v_{1}^{*}u^{*})^{*}\) ，我们有 \(uv_{1} \in \mathcal{L}_{1}(\mathrm{E}_{1};\mathrm{F})\) 且 \(\| uv_{1}\|_{1} \leqslant \| v_{1}^{*}\| \| u^{*}\|_{1} = \| v_{1}\| \| u\|_{1}\) （引理 8）。

命题由这些不等式立即得出。

命题 17. — 空间 \(\mathcal{L}_{1}(\mathrm{E};\mathrm{E})\) 与 E 的有限迹自同态所成的空间 \(\mathcal{L}_{1}(\mathrm{E})\) 相重合，且我们有 \(|\operatorname{Tr}(u)|\leqslant \| u\|_{1}\) ，对 E 的每个有限迹自同态 \(u\) 成立。

我们可以假设 K = C。按定义，空间 \(\mathcal{L}_{1}(E;E)\) 包含有限迹正自同态所成之集，因此 \(\mathcal{L}_{1}(E)\subset\mathcal{L}_{1}(E;E)\) （EVT, p. 50，定义 8）。反过来，我们要证明 \(\mathcal{L}_{1}(E;E)\) 含于 \(\mathcal{L}_{1}(E)\) 。

设 \(u \in \mathcal{L}_{1}(\mathrm{E};\mathrm{E})\) 。我们有 \(u^{*} \in \mathcal{L}_{1}(\mathrm{E};\mathrm{E})\) （引理 8）；因此只需证明 \(\mathcal{L}_{1}(\mathrm{E};\mathrm{E})\) 中的埃尔米特元素有有限迹（I, p. 96 的引理 2）。

设 u 是这样一个自同态。它是紧的（p. 169 的注 2）。设 B 是 E 的标准正交基，使得 u 在基 B 中是对角的（IV, p. 149 的定理 1），并设 \(\lambda\) 是 u 关于 B 的特征值族。自同态 \(|u|\) 在基 B 中可对角化，其特征值族是 \(|\lambda|\) （IV, p. 147 的命题 1, d)）。由于 \(|u|\) 有有限迹，后一族可和（IV, p. 164 的命题 12），因此族 \(\lambda\) 可和。于是 u 有有限迹（同前），且我们有 \(|\mathrm{Tr}(u)| \leqslant \mathrm{Tr}(|u|) = \|u\|_{1}\) 。

从而，空间 \(\mathcal{L}_{1}(\mathrm{E})\) 是 C*-代数 \(\mathcal{L}(\mathrm{E})\) 的自伴双边理想。若 E 是无限维的，这个理想在 \(\mathcal{L}(\mathrm{E})\) 中不闭。

命题 17 与命题 16 表明，映射 \(b\) ——从 \(\mathcal{L}_1(\mathrm{E};\mathrm{F})\times \mathcal{L}(\mathrm{F};\mathrm{E})\) 到 \(\mathbf{C}\) 中的映射——由 \(b(u,v) = \operatorname {Tr}(vu)\) 定义，是连续双线性型，它使空间 \(\mathcal{L}_1(\mathrm{E};\mathrm{F})\) 与 \(\mathcal{L}(\mathrm{F};\mathrm{E})\) 处于对偶。

引理 9. --- 设 \(u \in \mathcal{L}_1(\mathrm{E};\mathrm{F})\) 。我们有

\[
\| u\|_{1} = \sup_{\substack{v\in \mathcal{L}^{c}(\mathrm{F};\mathrm{E})\\ \| v\| \leqslant 1}}|b(u,v)|.
\]

由于每个希尔伯特--施密特映射都是紧的（IV, p. 165 的推论 2），我们有

\[
\| u\|_{1}\leqslant \sup_{\substack{v\in \mathcal{L}^{\mathrm{c}}(\mathrm{E};\mathrm{F})\\ \| v\| \leqslant 1}}|\operatorname{Tr}(vu)| = \sup_{\substack{v\in \mathcal{L}^{\mathrm{c}}(\mathrm{F};\mathrm{E})\\ \| v\| \leqslant 1}}|b(u,v)|
\]

（命题 14）。

设 \(v \in \mathcal{L}^{c}(\mathrm{E};\mathrm{F})\) 的范数 \(\leqslant 1\) 。由于 \(\mathcal{L}_{2}(\mathrm{E};\mathrm{F})\) 在 \(\mathcal{L}^{c}(\mathrm{E};\mathrm{F})\) 中稠密，存在序列 \((v_{n})_{n\in\mathbf{N}}\) （\(\mathcal{L}_{2}(\mathrm{E};\mathrm{F})\) 中的序列）在 \(\mathcal{L}(\mathrm{E};\mathrm{F})\) 中收敛到 v。序列 \((b(u,v_{n}))_{n\in\mathbf{N}}\) 收敛到 \(b(u,v)\) ；由于 \(|b(u,v_{n})| = |\operatorname{Tr}(v_{n}u)| = |\langle v_{n}^{*}|u\rangle| \leqslant \|u\|_{1}\) （命题 14）对所有 \(n \in N\) 成立，由此可得 \(|b(u,v)| \leqslant \|u\|_{1}\) 。

用 \(\theta\) 表示连续线性映射

\[
\theta \colon \mathcal {L} _ {1} (\mathrm{E}; \mathrm{F}) \to \mathcal {L} ^ {\mathrm{c}} (\mathrm{F}; \mathrm{E}) ^ {\prime}
\]

使得 \(\theta(u)(v) = b(u,v)\) 。

命题 18. --- 映射 θ 是等距同构。

由引理 9，映射 \(\theta\) 是等距的，只需证明 \(\theta\) 是满射。

设 \(\lambda \in \mathcal{L}^{\mathrm{c}}(\mathrm{F};\mathrm{E})'\) 。由于 \(\| v\| \leqslant \| v\|_{2}\) 对所有 \(v\in \mathcal{L}_2(\mathrm{F};\mathrm{E})\) 成立（EVT, V, p. 52，公式 (33)），\(\lambda\) 在子空间 \(\mathcal{L}_2(\mathrm{F};\mathrm{E})\) 上的限制是希尔伯特空间 \(\mathcal{L}_2(\mathrm{F};\mathrm{E})\) 上的连续线性型。因此存在元素 \(u\) （\(\mathcal{L}_2(\mathrm{F};\mathrm{E})\) 的元素），使得 \(\lambda (v) = \langle u|v\rangle\) 对所有 \(v\in \mathcal{L}_2(\mathrm{F};\mathrm{E})\) 成立（EVT, V, p. 15，定理 3）。

对每个 \(w \in \mathcal{L}_2(\mathrm{E};\mathrm{F})\) （范数 \(\leqslant 1\) 者），我们有 \(|\langle w|u\rangle| = |\lambda(w)| \leqslant \| \lambda \|\) 。于是 \(u\) 是从 \(\mathrm{E}\) 到 \(\mathrm{F}\) 中的核型映射（命题 14）。

连续线性型 \(\lambda\) 与 \(\theta(u)\) 在稠密子空间 \(\mathcal{L}_{2}(\mathrm{F},\mathrm{E})\) （\(\mathcal{L}^{\mathrm{c}}(\mathrm{F};\mathrm{E})\) 的稠密子空间）上相等。由此可得 \(\lambda=\theta(u)\) ，证明完成。

推论. --- 赋范空间 \(\mathcal{L}_{1}(\mathrm{E};\mathrm{F})\) 是巴拿赫空间。

由于空间 \(\mathcal{L}^{\mathrm{c}}(\mathrm{F};\mathrm{E})\) 是赋范空间，断言由本命题及 EVT, III, p. 24 的推论 2 得出。

\subsection{希尔伯特--施密特积分算子}

在本节中，X 与 Y 是局部紧拓扑空间。设 \(\mu\) 是 X 上的正测度，\(\nu\) 是 Y 上的正测度。我们分别用 \(\mathrm{L}^{2}(\mathrm{X})\) 、\(\mathrm{L}^{2}(\mathrm{Y})\) 与 \(\mathrm{L}^{2}(\mathrm{X} \times \mathrm{Y})\) 表示空间 \(\mathrm{L}^{2}(\mathrm{X}, \mu)\) 、\(\mathrm{L}^{2}(\mathrm{Y}, \nu)\) 与 \(\mathrm{L}^{2}(\mathrm{X} \times \mathrm{Y}, \mu \otimes \nu)\) ，对 \(\mathcal{L}^{2}(\mathrm{X})\) 、\(\mathcal{L}^{2}(\mathrm{Y})\) 与 \(\mathcal{L}^{2}(\mathrm{X} \times \mathrm{Y})\) 亦同样处理。

我们回顾（参见 III, p. 26 的 no. 4），对每个 \(N \in \mathcal{L}^{2}(X \times Y)\) ，存在唯一的连续线性映射 \(u_{\mathrm{N}}\) （从 \(L^{2}(X)\) 到 \(L^{2}(Y)\) 中的连续线性映射），使得

\[
\langle g \mid u _ {\mathrm{N}} (f) \rangle = \int_ {\mathrm{X} \times \mathrm{Y}} \overline {{g (y)}} \mathrm{N} (x, y) f (x) d (\mu \otimes \nu) (x, y)\tag{12}
\]

对所有 \(f \in \mathrm{L}^{2}(\mathrm{X})\) 与每个 \(g \in \mathrm{L}^{2}(\mathrm{Y})\) 成立。映射 \(u_{N}\) 是紧的（III, p. 33 的推论 1）。映射 \(N \mapsto u_{N}\) 过渡到商，诱导出从 \(\mathrm{L}^{2}(\mathrm{X} \times \mathrm{Y})\) 到 \(\mathcal{L}(\mathrm{L}^{2}(\mathrm{X});\mathrm{L}^{2}(\mathrm{Y}))\) 中的连续单射线性映射（III, p. 30 的命题 5）。

我们用 \(\theta\) 表示从 \(\mathcal{L}^2 (\mathrm{X})\otimes \mathcal{L}^2 (\mathrm{Y})\) 到 \(\mathcal{L}^2 (\mathrm{X}\times \mathrm{Y})\) 中的唯一线性映射，它对所有 \(f\in \mathcal{L}^2 (\mathrm{X})\) 与 \(g\in \mathcal{L}^2 (\mathrm{Y})\) 把 \(f\otimes g\) 映为由 \((x,y)\mapsto f(x)g(y)\) 定义的函数。

引理 10. --- 映射 \(\theta\) 过渡到商，定义了线性映射 \(\widetilde{\theta}\) （从 \(\mathrm{L}^2 (\mathrm{X})\otimes \mathrm{L}^2 (\mathrm{Y})\) 到 \(\mathrm{L}^2 (\mathrm{X}\times \mathrm{Y})\) 中的线性映射），且存在唯一的从 \(\mathrm{L}^2 (\mathrm{X})\widehat{\otimes}_2\mathrm{L}^2 (\mathrm{Y})\) 到 \(\mathrm{L}^2 (\mathrm{X}\times \mathrm{Y})\) 上的等距同构，它在 \(\mathrm{L}^2 (\mathrm{X})\otimes \mathrm{L}^2 (\mathrm{Y})\) 上与之相合。

第一个断言是初等的。我们来证明第二个。

设 \((f_i)_{i\in \mathrm{I}}\) 与 \((g_j)_{j\in \mathrm{J}}\) 分别是 \(\mathrm{L}^2 (\mathrm{X})\) 与 \(\mathrm{L}^2 (\mathrm{Y})\) 的标准正交基。族 \((\overline{f}_i\otimes g_j)_{(i,j)\in \mathrm{I}\times \mathrm{J}}\) 是 \(\mathrm{L}^2 (\mathrm{X})\widehat{\otimes}_2\mathrm{L}^2 (\mathrm{Y})\) 的标准正交基（EVT, V, p. 29，推论 1），而族 \((\widetilde{\theta} (\overline{f}_i\otimes g_j))_{(i,j)\in \mathrm{I}\times \mathrm{J}}\) 在 \(\mathrm{L}^2 (\mathrm{X}\times \mathrm{Y})\) 中标准正交。因此映射 \(\widetilde{\theta}\) 由连续性可扩张为从 \(\mathrm{L}^2 (\mathrm{X})\widehat{\otimes}_2\mathrm{L}^2 (\mathrm{Y})\) 到 \(\mathrm{L}^2 (\mathrm{X}\times \mathrm{Y})\) 中的等距线性映射。剩下要证明这个扩张是满射。

为此，只需证明 \(\widetilde{\theta}\) 的像在 \(\mathrm{L}^{2}(\mathrm{X}\times\mathrm{Y})\) 中稠密。设 N 是 \(\mathrm{L}^{2}(\mathrm{X}\times\mathrm{Y})\) 中与 \(\widetilde{\theta}\) 的像正交的元素。对所有 \(i\in I\) 与 \(j\in J\) ，我们有 \(\langle g_{j}\mid u_{\mathrm{N}}(f_{i})\rangle=\langle\widetilde{\theta}(\overline{f}_{i}\otimes g_{j})\mid\mathrm{N}\rangle=0\)。由公式 (12) 可得 \(u_{N}=0\) ，从而 N=0。

前一引理证实了 EVT, V, p. 29 例 2 中所述的断言。

在下文中，我们将借助引理 10 的等距同构把 \(\mathrm{L}^{2}(\mathrm{X})\widehat{\otimes}_{2}\mathrm{L}^{2}(\mathrm{Y})\) 与 \(\mathrm{L}^{2}(\mathrm{X}\times\mathrm{Y})\) 等同起来。

命题 19. --- 映射 \(N\mapsto u_{\mathrm{N}}\) 是从 \(\mathrm{L}^{2}(\mathrm{X}\times\mathrm{Y})\) 到空间 \(\mathcal{L}_{2}(\mathrm{L}^{2}(\mathrm{X});\mathrm{L}^{2}(\mathrm{Y}))\) （从 \(\mathrm{L}^{2}(\mathrm{X})\) 到 \(\mathrm{L}^{2}(\mathrm{Y})\) 中的希尔伯特--施密特映射所成的空间）上的等距同构。

从 \(\mathrm{L}^{2}(\mathrm{Y})\otimes\overline{\mathrm{L}^{2}(\mathrm{X})}\) 到 \(\mathcal{L}_{2}(\mathrm{L}^{2}(\mathrm{X});\mathrm{L}^{2}(\mathrm{Y}))\) 中把 \(g\otimes f\) 映为希尔伯特--施密特映射 \(h\mapsto\langle f|h\rangle g\) 的线性映射，可扩张为一个等距同构 \(\theta_{1}\) ，即从 \(\mathrm{L}^{2}(\mathrm{Y})\widehat{\otimes}_{2}\overline{\mathrm{L}^{2}(\mathrm{X})}\) 到 \(\mathcal{L}_{2}(\mathrm{L}^{2}(\mathrm{X});\mathrm{L}^{2}(\mathrm{Y}))\) 上的等距同构（EVT, V, p. 52，定理 1）。

另外，用 \(\theta_{2}\) 表示 \(\mathrm{L}^{2}(\mathrm{X})\widehat{\otimes}_{2}\mathrm{L}^{2}(\mathrm{Y})\) 到 \(\mathrm{L}^{2}(\mathrm{Y})\widehat{\otimes}_{2}\overline{\mathrm{L}^{2}(\mathrm{X})}\) 上的等距同构，它把 \(f\otimes g\) 映为 \(g\otimes \overline{f}\) ，对所有 \(f\in \mathrm{L}^2 (\mathrm{X})\) 与每个 \(g\in \mathrm{L}^2 (\mathrm{Y})\) 成立。

线性映射 \(\theta_{3} = \theta_{1} \circ \theta_{2}\) 等同于 \(\mathrm{L}^{2}(\mathrm{X} \times \mathrm{Y})\) 到 \(\mathcal{L}_{2}(\mathrm{L}^{2}(\mathrm{X}); \mathrm{L}^{2}(\mathrm{Y}))\) 上的一个等距同构。

设 \(f \in \mathrm{L}^2(\mathrm{X})\) 且 \(g \in \mathrm{L}^2(\mathrm{Y})\) ，并设 \(\mathrm{N}\) 是 \(\mathrm{L}^2(\mathrm{X} \times \mathrm{Y})\) 的元素（与 \(f \otimes g\) 等同）。由 INT, V, p. 95, § 8, no. 3 的推论 2 及公式 (12)，我们有

\[
\langle g _ {1} \mid u _ {\mathrm{N}} (f _ {1}) \rangle = \langle \overline {{f}} \mid f _ {1} \rangle \langle g _ {1} \mid g \rangle
\]

对所有 \(f_1 \in \mathrm{L}^2(\mathrm{X})\) 与 \(g_1 \in \mathrm{L}^2(\mathrm{Y})\) 成立。映射 \(u = \theta_3(\mathrm{N})\) 就是线性映射 \(\theta_1(g \otimes \overline{f})\) ；因此它满足 \(u(h) = \langle \overline{f} | h \rangle g\) ，对所有 \(h \in \mathrm{L}^2(\mathrm{X})\) 成立。于是，对所有 \(f_1 \in \mathrm{L}^2(\mathrm{X})\) 与 \(g_1 \in \mathrm{L}^2(\mathrm{Y})\) ，可得

\[
\langle g _ {1} \mid u (h _ {1}) \rangle = \langle \overline {{f}} \mid h _ {1} \rangle \langle g _ {1} \mid g \rangle = \langle g _ {1} \mid u _ {\mathrm{N}} (h _ {1}) \rangle ,
\]

由此得 \(\theta_3(N) = u_N\) 。由于 \(\theta_3\) 与 \(N \mapsto u_N\) 都连续，我们断定 \(\theta_3(N) = u_N\) 对所有 \(N \in L^2(X \times Y)\) 成立，证明完成。

推论. --- 对每个 \(N \in L^{2}(X \times Y)\) ，线性映射 \(u_{N}\) 是希尔伯特--施密特算子，且我们有 \(\mathrm{Tr}(u_{\mathrm{N}}^{*}u_{\mathrm{N}}) = \|N\|^{2}\) 。

事实上，\(\| u\| _2 = \sqrt{\mathrm{Tr}(u^*u)}\) 对所有 \(u\in \mathcal{L}_2(\mathrm{L}^2 (\mathrm{X});\mathrm{L}^2 (\mathrm{Y}))\) 成立。

注. --- 设 \(N \in L^{2}(X \times Y)\) 。由 IV, p. 150 的推论 2，存在可数集 I、标准正交族 \((f_{i})_{i \in I}\) （\(L^{2}(X)\) 中的）与标准正交族 \((g_{i})_{i \in I}\) （\(L^{2}(Y)\) 中的），以及族 \((\alpha_{i})_{i \in I}\) （\(\mathbf{R}_{+}^{*}\) 中的族），使得

\[
u _ {\mathrm{N}} (f) = \sum_ {i \in \mathrm{I}} \alpha_ {i} \langle f _ {i} | f \rangle g _ {i}
\]

对所有 \(f \in \mathrm{L}^2(\mathrm{X})\) 成立，其中级数在 \(\mathrm{L}^2(\mathrm{Y})\) 中收敛。由 IV, p. 165 的推论 4 及上述推论，我们有

\[
\sum_ {i \in \mathrm{I}} \alpha_ {i} ^ {2} = \mathrm{Tr} (u _ {\mathrm{N}} ^ {*} u _ {\mathrm{N}}) = \| u _ {\mathrm{N}} \| _ {2} ^ {2} = \int_ {\mathrm{X} \times \mathrm{Y}} | \mathrm{N} (x, y) | ^ {2} d (\mu \otimes \nu) (x, y).
\]

另外，注意 \(h_{i,j} = \overline{f}_{i} \otimes g_{j} \in \mathrm{L}^{2}(\mathrm{X} \times \mathrm{Y})\) 对所有 \((i, j) \in \mathrm{I} \times \mathrm{I}\) 成立。于是我们有

\[
\mathrm{N} = \sum_ {i \in \mathrm{I}} \alpha_ {i} h _ {i, i}
\]

in \(\mathrm{L}^2 (\mathrm{X}\times \mathrm{Y})\) .

事实上，设 \((f_{j})_{j\in\mathrm{J}}\) 与 \((g_{k})_{k\in\mathrm{K}}\) 分别是 \(\mathrm{L}^{2}(\mathrm{X})\) 与 \(\mathrm{L}^{2}(\mathrm{Y})\) 的标准正交基，且分别是族 \((f_{i})_{i\in\mathrm{I}}\) 与 \((g_{i})_{i\in\mathrm{I}}\) 的扩张。置 \(h_{j,k}=\overline{f}_{j}\otimes g_{k}\) ，对所有 \((j,k)\in\mathrm{J}\times\mathrm{K}\) 。由引理 10，族 \((h_{j,k})_{(j,k)\in\mathrm{J}\times\mathrm{K}}\) 是 \(\mathrm{L}^{2}(\mathrm{X}\times\mathrm{Y})\) 的标准正交基。我们有 \(\langle h_{j,k}\mid\mathrm{N}\rangle=\langle g_{k}\mid u_{\mathrm{N}}(f_{j})\rangle\) ，对所有 \((j,k)\in\mathrm{J}\times\mathrm{K}\) 。若 \(j\notin I\) ，这个量为零。若 \(j\in I\) ，它等于 \(\alpha_{j}\langle g_{k}\mid g_{j}\rangle\) ，故除 k=j 外为零，而当 k=j 时 \(\langle h_{j,j}\mid N\rangle=\alpha_{j}\) 。因此

\[
\mathrm{N} = \sum_ {(j, k) \in \mathrm{J} \times \mathrm{K}} \left\langle h _ {j, k} \mid \mathrm{N} \right\rangle h _ {j, k} = \sum_ {i \in \mathrm{I}} \alpha_ {i} h _ {i, i}.
\]

\subsection{连续核积分算子的迹}

在本节中，我们保持前一节的约定，取 Y = X 且 \(\nu = \mu\)。我们假设 X 是在无穷远处可数的局部紧拓扑空间（TG, I, p. 68，定义 5）。

特别地，我们把空间 \(\mathrm{L}^{2}(\mathrm{X}\times\mathrm{X})\) 与 \(\mathrm{L}^{2}(\mathrm{X})\widehat{\otimes}_{2}\mathrm{L}^{2}(\mathrm{X})\) 等同起来（IV, p. 172 的引理 10）。我们用 \(\widetilde{f}\) 表示 \(\mathrm{L}^{2}(\mathrm{X})\) （相应地，\(\mathrm{L}^{2}(\mathrm{X}\times\mathrm{X})\) ）中函数 \(f\in\mathcal{L}^{2}(\mathrm{X})\) （相应地，\(\mathcal{L}^{2}(\mathrm{X}\times\mathrm{X})\) 中的函数）的类。

命题 20. --- 设 \((f_n)_{n \in \mathbf{N}}\) 是从 X 到可度量化空间 F 中的可测映射序列。假设 \(f_n(x)\) 的极限在 X 的一个 \(\mu\)-零测子集之外存在。则存在 X 的紧子集序列 \((\mathrm{C}_m)_{m \in \mathbf{N}}\) ，其并 \(\widetilde{\mathrm{X}}\) 满足 \(|\mu|(\mathrm{X} - \widetilde{\mathrm{X}}) = 0\) ，使得函数 \(f_n\) 在 \(\mathrm{C}_m\) 上连续（对所有 \(n \in \mathbf{N}\) 与每个 \(m \in \mathbf{N}\)），且序列 \((f_n)_{n \in \mathbf{N}}\) 一致收敛到 \(f\) （在 \(\mathrm{C}_m\) 上，对所有 \(m \in \mathbf{N}\)）。

由 INT, IV, p. 175, § 5, no. 4 的定理 2 及 INT, IV, p. 188, § 5, no. 8 的命题 12, b)，X 中使得函数 \(f_n\) 对每个 \(n\) 在 C 上连续且 \((f_n)\) 在 C 上一致收敛到 \(f\) 的紧子集 C 所成之集在 X 中 \(\mu\)-稠密（INT, IV, p. 189, § 5, no. 8，定义 6）。断言随即由 INT, IV, p. 189, § 5, no. 8 的注得出。

设 \(N\) 是属于 \(\mathcal{L}^2 (\mathrm{X}\times \mathrm{X})\) 的函数，\(u_{\mathrm{N}}\) 是从 \(\mathrm{L}^2 (\mathrm{X})\) 到 \(\mathrm{L}^2 (\mathrm{X})\) 中的以 \(N\) 为核的希尔伯特--施密特映射（IV, p. 173 的命题 19）。由于映射 \(u_{\mathrm{N}}\) 是紧的，由 IV, p. 156 的命题 9，存在元素 \(M\) （\(\overline{\mathbf{N}}\) 中的元素），并且用 \(I\) 表示整数 \(n\in \mathbf{N}\) 中满足 \(n\leqslant \mathrm{M}\) 者所成之集，存在标准正交族 \((f_i)_{i\in \mathrm{I}}\) 与 \((g_{i})_{i\in \mathrm{I}}\) （\(\mathcal{L}^2 (\mathrm{X})\) 中的标准正交族）以及族 \((\alpha_{i})_{i\in \mathrm{I}}\) （\(\mathbf{R}_+^*\) 中的族），使得

\[
u _ {\mathrm{N}} (f) = \sum_ {i \in \mathrm{I}} \alpha_ {i} \langle \widetilde {f} _ {i} | f \rangle \widetilde {g} _ {i}\tag{13}
\]

对所有 \(f \in \mathrm{L}^2(\mathrm{X})\) 成立，其中级数在 \(\mathrm{L}^2(\mathrm{X})\) 中收敛。

对每个 \(i \in I\) ，用 \(h_i \in \mathcal{L}^2(\mathrm{X} \times \mathrm{X})\) 表示由 \(h_i(x, y) = \overline{f_i(x)} g_i(y)\) 定义的函数，从而 \(\widetilde{h}_i\) 是 \(\overline{f}_i \otimes g_i\) 的类。由 IV, p. 173 的注 9，通项为 \(\alpha_i h_i\) 的级数在 \(\mathrm{L}^2(\mathrm{X} \times \mathrm{X})\) 中收敛到 N。

在本节其余部分，我们假设 N 是连续函数。

命题 21. --- 假设 \(u_{\mathrm{N}}\) 有有限迹。则存在一个集 \(\widetilde{\mathrm{X}} \subset \mathrm{X}\) ，其余集 \(\mathrm{X} - \widetilde{\mathrm{X}}\) 是 \(\mu\)-零测的，还存在函数 \(\mathrm{H} \in \mathcal{L}^2(\mathrm{X} \times \mathrm{X})\) ，满足以下条件：

\begin{enumerate}
\def\labelenumi{(\roman{enumi})}
\item
  对每个 \((x,y)\in\widetilde{\mathrm{X}}\times\widetilde{\mathrm{X}}\) ，族 \((\alpha_{i}\overline{f_{i}(x)}g_{i}(y))_{i\in\mathrm{I}}\) 在 \(\mathbf{C}\) 中可和，其和为 \(\mathrm{N}(x,y)\) ；
\item
  对 I 的每个有限子集 J 与每个 \((x,y)\in\widetilde{\mathbf{X}}\times\widetilde{\mathbf{X}}\) ，我们有
\end{enumerate}

\[
\left| \sum_ {i \in \mathrm{J}} \alpha_ {i} h _ {i} (x, y) \right| \leqslant \mathrm{H} (x, y);\tag{14}
\]

\begin{enumerate}
\def\labelenumi{(\roman{enumi})}
\setcounter{enumi}{2}
\tightlist
\item
  函数 \(x \mapsto \mathrm{H}(x, x)\) 属于 \(\mathcal{L}^1(\mathrm{X})\) 。
\end{enumerate}

由于 \(u_{N}\) 有有限迹，族 \((\alpha_{i})\) 可和（IV, p. 165 的推论 4）。级数

\[
\sum_ {i \in \mathrm{I}} \alpha_ {i} | f _ {i} | ^ {2}, \quad \sum_ {i \in \mathrm{I}} \alpha_ {i} | g _ {i} | ^ {2}
\]

因此在 \(\mathcal{L}^{1}(X)\) 中收敛，从而 \(\mu\)-几乎处处收敛。用 F 与 G 分别表示由这些级数的和在 \(\mu\)-几乎处处意义下定义的函数，其中置 \(F(x) = 0\) 与 \(G(x) = 0\) ，对相应级数不收敛的每个 \(x\) 。由于 F 与 G 属于 \(\mathcal{L}^{1}(X)\)，由 \(H(x,y) = \sqrt{F(x)G(y)}\) 定义的函数 H 属于 \(\mathcal{L}^{2}(X \times X)\) （INT, V, p. 95, § 8, no. 3，推论 2）。

把命题 20 应用于空间 X，取 \(F = \mathbf{C}^{2}\) ，并取可测映射

\[
s_{n}\colon x\mapsto \Bigl (\sum_{\substack{i\in \mathrm{I}\\ i\leqslant n}}\alpha_{i}|f_{i}(x)|^{2},\sum_{\substack{i\in \mathrm{I}\\ i\leqslant n}}\alpha_{i}|g_{i}(x)|^{2}\Bigr)
\]

它们定义在 X 上。因此存在 X 的紧子集序列 \((\mathrm{C}_{m})_{m\in\mathbf{N}}\) ，满足以下条件：

\begin{enumerate}
\def\labelenumi{(\arabic{enumi})}
\item
  并 \(\widetilde{X}\) 满足 \(\mu(\mathrm{X}-\widetilde{\mathrm{X}})=0\) ；
\item
  对所有 \(m \in \mathbf{N}\) 与每个 \(n \in \mathbf{N}\) ，函数 \(s_n\) 在 \(\mathrm{C}_m\) 上连续；
\item
  对所有 \(m \in \mathbf{N}\) ，级数 \(\sum \alpha_{i}|f_{i}|^{2}\) 与 \(\sum \alpha_{i}|g_{i}|^{2}\) 在 \(C_{m}\) 上分别一致收敛到 F 与 G；
\item
  由 \(\mu\) 在 \(C_{m}\) 上诱导的测度的支撑等于 \(C_{m}\) （必要时把 \(C_{m}\) 换成这个测度的支撑）。
\end{enumerate}

我们来证明集 \(\widetilde{\mathbf{X}}\) 与函数 H 满足条件 (i)、(ii) 与 (iii)。

设 \((x,y)\in \widetilde{\mathbf{X}}\times \widetilde{\mathbf{X}}\)。对每个有限子集 \(\mathrm{J}\subset \mathrm{I}\) ，我们有

\[
\left| \sum_ {i \in \mathrm{J}} \alpha_ {i} \overline {{f _ {i} (x)}} g _ {i} (y) \right| \leqslant \left(\sum_ {i \in \mathrm{J}} \alpha_ {i} | f _ {i} (x) | ^ {2}\right) ^ {1 / 2} \left(\sum_ {i \in \mathrm{J}} \alpha_ {i} | g _ {i} (y) | ^ {2}\right) ^ {1 / 2}\tag{15}
\]

进而有

\[
\left| \sum_ {i \in \mathrm{J}} \alpha_ {i} \overline {{f _ {i} (x)}} g _ {i} (y) \right| \leqslant \mathrm{H} (x, y),\tag{16}
\]

这已给出性质 (ii)。

由不等式 (15)，级数 \(\sum_{i} \alpha_{i} h_{i}\) 在 \(\mathrm{K}_{m} \times \mathrm{K}_{n}\) 上一致收敛（对所有 \((m, n) \in \mathbf{N}^{2}\)）。设 \(\widetilde{\mathrm{N}}\) 是 \(\mathrm{X} \times \mathrm{X}\) 上由

\[
\widetilde {\mathrm{N}} (x, y) = \sum_ {i \in \mathrm{I}} \alpha_ {i} h _ {i} (x, y) = \sum_ {i \in \mathrm{I}} \alpha_ {i} \overline {{f _ {i} (x)}} g _ {i} (y)
\]

对所有 \((x,y)\in\widetilde{\mathrm{X}}\times\widetilde{\mathrm{X}}\) 定义、否则置 \(\widetilde{\mathrm{N}}(x,y)=0\) 的函数。函数 \(\widetilde{N}\) 可测（INT, IV, p. 175, § 5, n°4，定理 2），且在 \(K_{m}\times K_{n}\) 上连续（对所有 \((m,n)\in\mathbf{N}^{2}\)）。

由 (16)，我们有 \(|\widetilde{\mathrm{N}}(x,y)| \leqslant \mathrm{H}(x,y)\) ，对所有 \((x,y) \in \mathrm{X} \times \mathrm{X}\) 成立，特别地 \(\widetilde{\mathrm{N}}\) 属于 \(\mathcal{L}^2(\mathrm{X} \times \mathrm{X})\) （INT, IV, p. 84, § 5, no. 6，定理 5）。

我们来证明 \(N = \widetilde{N}\) 在 \(\widetilde{X} \times \widetilde{X}\) 上成立，这将给出性质 (i)。设 f 与 g 是 \(\mathcal{L}^{2}(X)\) 的元素。我们有

\[
\langle g \mid u _ {\widetilde {N}} (f) \rangle = \int_ {X \times X} \widetilde {N} (x, y) f (x) \overline {{g (y)}} d (\mu \otimes \mu) (x, y)
\]

（IV, p. 172 的公式 (12)）。对每个 \((x,y)\in \widetilde{\mathrm{X}}\times \widetilde{\mathrm{X}}\) 与 I 的每个有限子集 J，我们有

\[
\left| \sum_ {i \in J} \alpha_ {i} \overline {{f _ {i} (x)}} g _ {i} (y) f (x) \overline {{g (y)}} \right| \leqslant | f (x) g (y) | H (x, y)
\]

（公式 (16)）。由于这个不等式的右端在 X × X 上可积（INT, V, p. 95, § 8, n° 3，推论 2），可以应用勒贝格定理（INT, IV, p. 137, § 3, n° 7，定理 6）及公式 (13)，推出

\[
\begin{array}{c} \langle g | u _ {\widetilde {N}} (f) \rangle = \sum_ {i \in I} \alpha_ {i} \int_ {X \times X} \overline {{f _ {i} (x)}} g _ {i} (y) f (x) \overline {{g (y)}} d (\mu \otimes \mu) (x, y) \\ = \sum_ {i \in I} \alpha_ {i} \langle f _ {i} | f \rangle \langle g | g _ {i} \rangle = \langle g | u _ {N} (f) \rangle . \end{array}
\]

因此可得 \(u_{\widetilde{\mathrm{N}}} = u_{\mathrm{N}}\) ，从而 \(\mathrm{N} = \widetilde{\mathrm{N}}\) 在 \(\mathrm{L}^2(\mathrm{X} \times \mathrm{X})\) 中成立 （III, p. 28 的命题 3, b)））。

对每个 \((m,n)\in\mathbf{N}^{2}\)，函数 N 与 \(\widetilde{N}\) 在 \(K_{m}\times K_{n}\) 上连续，故在 \(K_{m}\times K_{n}\) 上相等（INT, III, p. 69, § 2, n°2 的命题 9），因为 \(\mu\otimes\mu\) 在 \(K_{m}\times K_{n}\) 上诱导的测度的支撑等于 \(K_{m}\times K_{n}\) 。因此 N 与 \(\widetilde{N}\) 在 \(\widetilde{X}\times\widetilde{X}\) 上相合。

最后，上界 \(\sqrt{\mathrm{FG}} \leqslant (\mathrm{F} + \mathrm{G}) / 2\) 蕴含函数 \(x \mapsto \sqrt{\mathrm{F}(x)\mathrm{G}(x)} = \mathrm{H}(x,x)\) 在 X 上可积，由此得到性质 (iii)。

在本节其余部分，我们保持该命题的记号。

定理 2. --- 假设 \(u_{\mathrm{N}}\) 有有限迹。则函数 \(x \mapsto \mathrm{N}(x, x)\) 属于 \(\mathcal{L}^1(\mathrm{X})\) ，且我们有

\[
\mathrm{Tr} (u _ {\mathrm{N}}) = \int_ {\mathrm{X}} \mathrm{N} (x, x) d \mu (x).
\]

由条件 (i) 与 (ii)，我们有 \(|\mathrm{N}(x,x)| \leqslant \mathrm{H}(x,x)\) （对 \(x \in \widetilde{X}\)），故由条件 (iii)，函数 \(x \mapsto \mathrm{N}(x,x)\) 属于 \(\mathcal{L}^{1}(\mathrm{X})\) 。

条件 (i) 给出等式

\[
\mathrm{N} (x, x) = \sum_ {i \in \mathrm{I}} \alpha_ {i} \overline {{f _ {i} (x)}} g _ {i} (x)
\]

对所有 \(x \in \widetilde{\mathrm{X}}\) 成立。由条件 (ii) 与 (iii)，可以应用勒贝格定理（INT, IV, p. 137, § 3, no. 7，定理 6），由此可得

\[
\begin{array}{r} \int_ {\mathrm{X}} \mathrm{N} (x, x) d \mu (x) = \sum_ {i \in \mathrm{I}} \alpha_ {i} \int_ {\mathrm{X}} \overline {{f _ {i} (x)}} g _ {i} (x) d \mu (x) \\ = \sum_ {i \in \mathrm{I}} \alpha_ {i} \langle f _ {i} | g _ {i} \rangle = \mathrm{Tr} (u _ {\mathrm{N}}), \end{array}
\]

最后一个等式来自 IV, p. 170 命题 17 的推论。

引理 11. --- 若自同态 \(u_{\mathrm{N}}\) （\(\mathrm{L}^2(\mathrm{X})\) 的自同态）是正的，则 \(\mathrm{N}(x,x) \geqslant 0\) 对每个 \(x\) （\(\mu\) 的支撑中的元素）成立。

由于 \(u_{N}\) 是正的，可以选取族 \((f_{i})\) 与 \((g_{i})\) ，使得 \(f_{i} = g_{i}\) 对所有 \(i \in I\) 成立（参见 IV, p. 151 的注 3）。于是对每个 \(x \in \widetilde{X}\) ，我们有

\[
\mathrm{N} (x, x) = \sum_ {i \in \mathrm{I}} \alpha_ {i} | f _ {i} (x) | ^ {2} \geqslant 0.
\]

由于 \(N\) 连续且 \(\mu(\mathbf{X}-\widetilde{\mathbf{X}})=0\) ，函数 \(x \mapsto N(x,x)\) 在 \(\mu\) 的支撑上非负。

注. --- 该引理的逆命题不成立（IV, p. 316 的习题 14）。

命题 22. --- 假设 \(u_{\mathrm{N}}\) 是 \(\mathrm{L}^2(\mathrm{X})\) 的正自同态。则自同态 \(u_{\mathrm{N}}\) 有有限迹当且仅当函数 \(x \mapsto \mathrm{N}(x,x)\) 是 \(\mu\)-可积的。此时我们有

\[
\mathrm{Tr} (u _ {\mathrm{N}}) = \int_ {\mathrm{X}} \mathrm{N} (x, x) d \mu (x).
\]

若 \(u_{N}\) 有有限迹，定理 2 蕴含 \(x \mapsto \mathrm{N}(x, x)\) 在 X 上可积，且其积分是 \(u_{N}\) 的迹。

反之，设函数 \(x \mapsto \mathrm{N}(x, x)\) 可积。由于 \(u_{N}\) 是正的，可以选取标准正交族 \((f_i)_{i \in I}\) 与 \((g_i)_{i \in I}\) ，使得 \(f_i = g_i\) 对所有 \(i \in I\) 成立（IV, p. 151 的注 3）。对每个 \(x \in \widetilde{X}\) ，我们有

\[
\mathrm{N} (x, x) = \sum_ {i \in \mathrm{I}} \alpha_ {i} | f _ {i} (x) | ^ {2},
\]

由此，对 I 的每个有限子集 J，我们推出

\[
\begin{array}{c} \sum_ {i \in \mathrm{J}} \alpha_ {i} = \sum_ {i \in \mathrm{J}} \alpha_ {i} \int_ {\mathrm{X}} | f _ {i} (x) | ^ {2} d \mu (x) \\ = \int_ {\mathrm{X}} \sum_ {i \in \mathrm{J}} \alpha_ {i} | f _ {i} (x) | ^ {2} d \mu (x) \leqslant \int_ {\mathrm{X}} \mathrm{N} (x, x) d \mu (x). \end{array}
\]

因此族 \((\alpha_{i})_{i\in\mathrm{I}}\) 可和，这蕴含自同态 \(u_{N}\) 有有限迹（IV, p. 164 的命题 12）。

注. --- 即使当 X 紧且 N 连续时，自同态 \(u_{\mathrm{N}}\) （\(\mathrm{L}^{2}(\mathrm{X})\) 的自同态）也不总有有限迹（参见 IV, p. 314 的习题 8）。

\section{正规自同态}

在本章中，除另有说明外，所考虑的希尔伯特空间都是复的。

\subsection{\texorpdfstring{关于空间 \(L^{p}(X,\mu)\) 的补充}{关于空间 L\^{}\{p\}(X,\textbackslash mu) 的补充}}

命题 1. --- 设 X 是局部紧拓扑空间。设 \(\mu\) 是 X 上的正测度，\(p \in [1, +\infty[\)。设 \((\mathrm{X}_i)_{i \in \mathrm{I}}\) 是 X 中两两不相交的局部紧子集的局部可数族（INT, IV, p. 190, § 5, no. 9，定义 7），使得诸 \(\mathrm{X}_i\) 之并的余集是局部 \(\mu\)-零测的。设 \(\varphi_i\) 是 \(\mathrm{X}_i\) 的特征函数，\(\mu_i\) 是 \(\mu\) 在 \(\mathrm{X}_i\) 上诱导的测度（INT, IV, p. 186, §5, n° 7，定义 4）。

\begin{enumerate}
\def\labelenumi{\alph{enumi})}
\item
  映射 \(p_{i}\) ：\(f \mapsto f\varphi_{i}\) 是 \(\mathcal{L}^{p}(\mathrm{X}, \mu)\) 的投影子，且过渡到商，定义了投影子 \(\widetilde{p}_{i}\) （\(\mathrm{L}^{p}(\mathrm{X}, \mu)\) 的投影子）。当 p = 2 时，后者是 \(\mathrm{L}^{2}(\mathrm{X}, \mu)\) 的正交投影子；
\item
  映射 \(f \mapsto f|X_i\) 诱导出从 \(p_i\) 的像到空间 \(\mathcal{L}^p(X_i, \mu_i)\) 上的等距同构，且过渡到商，定义了从 \(\widetilde{p}_i\) 的像到 \(\mathrm{L}^p(X_i, \mu_i)\) 上的等距同构，借助它这两个空间互相等同；
\item
  诸空间 \(\mathrm{L}^{p}(\mathrm{X}_{i},\mu_{i})\) 之和在空间 \(\mathrm{L}^{p}(\mathrm{X},\mu)\) 中是完全的。当 p = 2 时，空间 \(\mathrm{L}^{2}(\mathrm{X},\mu)\) 是诸空间 \(\mathrm{L}^{2}(\mathrm{X}_{i},\mu_{i})\) 的希尔伯特和。
\end{enumerate}

断言 \(a\)）是初等的。断言 \(b\)）由 INT, V, p. 84, § 7, n° 1 的评注得出。

设 \(q \in ]1, +\infty]\) 是 \(p\) 的共轭指数。我们把 \(L^p(X, \mu)\) 的对偶与 \(L^q(X, \mu)\) 等同起来（INT, V, p. 61, § 5, n° 8，定理 4）。设 \(f\) 是 \(\mathcal{L}^q(X, \mu)\) 中的函数，其类 \(\widetilde{f}\) 在 \(L^q(X, \mu)\) 中满足 \(\langle \widetilde{f}, \widetilde{p}_i(\varphi) \rangle = 0\) （对所有 \(i \in I\) 与每个 \(\varphi \in L^p(X, \mu)\)）。由于 \(p_i\) 的像包含 \(\mathcal{K}(X_i)\) ，可得测度 \((f|X_i) \cdot \mu_i\) 在 \(X_i\) 上为零，即 \(f\) 在 \(X_i\) 上的限制关于 \(\mu_i\) 在 \(X_i\) 上局部零测（INT, V, p. 46, § 5, n° 3，推论 2）。由于诸 \(X_i\) 之并在 X 中的余集是局部 \(\mu\)-零测的，可以断定函数 \(f\) 在 X 上局部 \(\mu\)-零测（INT, IV, p. 190, § 5, n° 9 及 p. 172, n° 2，命题 5）。于是 \(f\) 的类在 \(L^q(X, \mu)\) 中为零（当 \(p \neq 1\) 时用到 INT, V, p. 8, § 1, n° 3 的引理 1 及命题 9 的推论）。由哈恩--巴拿赫定理（EVT, II, p. 46，推论 1），断言 \(c\)）的第一部分得证。

若 \(p = 2\) ，\(p_i\) 的像与 \(p_j\) 的像正交（对所有 \(i \neq j\) ），因为此时 \(X_i\) 与 \(X_j\) 不相交，由此得到最后一个断言。

命题 2. --- 设 X 是在无穷远处可数的局部紧空间。设 \(\mu\) 是 X 上的正测度。对每个 \(p \in [1, +\infty[\)，空间 \(\mathrm{L}^{p}(\mathrm{X},\mu)\) 是可数型的。

设 \((\mathrm{U}_{n})_{n\in\mathbf{N}}\) 是 X 中相对紧开集的序列，其并等于 X 且满足 \(\overline{U}_{n}\subset U_{n+1}\) 对所有 \(n\in\mathbf{N}\) 成立（TG, I, p. 68，命题 15）。对每个 \(n\in\mathbf{N}\)，空间 \(\mathcal{K}(\mathrm{X},\overline{\mathrm{U}}_{n})\) 等同于巴拿赫空间 \(\mathcal{C}(\overline{\mathrm{U}}_{n})\) 的闭子空间（INT, III, p. 40, § 1, n° 1）；由于后一空间是可数型的（TG, X, p. 25，推论），\(\mathcal{K}(\mathrm{X},\overline{\mathrm{U}}_{n})\) 亦然（TG, IX, p. 19, cor., (i)）。设 \(F_{n}\) 是 \(\mathcal{K}(\mathrm{X},\overline{\mathrm{U}}_{n})\) 的可数稠密子集。

设 \(f \in \mathcal{L}^p(\mathrm{X},\mu)\) 且 \(\varepsilon > 0\) 。存在整数 \(n \in \mathbf{N}\) 使得 \(\int_{\mathrm{X}-\mathrm{U}_n}|f|^p < \varepsilon / 2\) ，且存在 \(g \in \mathcal{F}_n\) 使得 \(\int_{\overline{\mathrm{U}}_n}|f - g|^p < \varepsilon / 2\) 。因此，\(\mathrm{L}^p(\mathrm{X},\mu)\) 中诸集合 \(\mathcal{F}_n\) 的元素之类所成的并在 \(\mathrm{L}^p(\mathrm{X},\mu)\) 中稠密，证明完成（TG, IX, p. 18，命题 12）。

\subsection{可测函数的本质像}

在本节中，我们用 X 表示局部紧拓扑空间，用 \(\mu\) 表示 X 上的正测度。

定义 1. --- 设 Y 是拓扑空间。对每个从 X 到 Y 中的可测函数 g，\(\mu\)-本质像（\(g\) 的）是满足如下条件的 \(y \in Y\) 所成之子集：对每个开邻域 \(U\) （\(y\) 的），集 \(g^{-1}(\mathrm{U})\) 在 X 中不是局部 \(\mu\)-零测的（INT, IV, p. 172, § 5, no. 2，定义 3）。

引理 1. --- 设 \(g\) 是从 X 到拓扑空间 Y 中的 \(\mu\)-可测函数，并设 \(S\) 是它的 \(\mu\)-本质像。元素 \(x \in X\) 中使得 \(g(x)\) 不属于 S 者所成之集在 X 中局部 \(\mu\)-零测。

设 \(Z = g^{-1}(Y - S)\) 是所考虑的集。

首先假设 X 紧且 g 连续。在此情形，前推测度 \(g(\mu)\) 有定义（INT, V, p. 69, § 6, no. 1，定义 1），因为 \(\mu\) 是有界测度。由定义可得，\(\mu\)-本质像（\(g\) 的）是测度 \(g(\mu)\) 的支撑（参见 INT, V, p. 70, § 6, no. 2，推论 2），于是 \(\mu(Z) = g(\mu)(Y - S) = 0\) （INT, IV, p. 118, § 2, no. 2，命题 5）。

现在考虑一般情形。设 C 是 X 的紧子集，\(\varepsilon > 0\) 是实数。由于 \(g\) 可测，存在紧子集 \(C_1 \subset C\) ，使得 \(\mu(C - C_1) \leqslant \varepsilon\) 且 \(g|C_1\) 连续（INT, IV, p. 169, § 5, no. 1，命题 1）。于是我们有

\[
\mu (\mathrm{Z} \cap \mathrm{C}) \leqslant \mu (\mathrm{C} - \mathrm{C} _ {1}) + \mu (\mathrm{Z} \cap \mathrm{C} _ {1}) \leqslant \varepsilon + \mu (\mathrm{Z} \cap \mathrm{C} _ {1}).
\]

设 \(\mu_{1}\) 是 \(\mu\) 在 \(C_{1}\) 上诱导的测度（INT, IV, p. 186, § 5, no. 7，定义 4）。设 \(S_{1}\) 是 \(\mu_{1}\)-本质像（\(g|C_{1}\) 的）。我们有 \(Z\cap C_{1}\subset(g|C_{1})^{-1}(Y-S_{1})\) 。由第一种情形，集 \(Z\cap C_{1}\) 因此是 \(\mu_{1}\)-零测的，从而是 \(\mu\)-零测的（INT, IV, p. 187, § 5, no. 7，引理 2, (i)）。上面的不等式变成 \(\mu(Z\cap C)\leqslant\varepsilon\) ；由于 \(\varepsilon\) 与 C 是任意的，集 Z 局部 \(\mu\)-零测（INT, IV, p. 172, § 5, no. 2，命题 5）。

引理 2. --- 设 \(g\) 是从 X 到 \(\mathbf{C}\) 中的连续函数。\(\mu\)-本质像（\(g\) 的）是 \(\mu\) 的支撑在 \(g\) 下的像的闭包。

设 \(Y\) 是 \(\mu\) 的支撑。若 \(z \in \mathbf{C}\) 不附着于 \(g(Y)\)，则存在开邻域 \(U\) （\(z\) 的），使得开集 \(g^{-1}(U)\) 不与 Y 相交，从而是局部 \(\mu\)-零测的；这意味着 \(z\) 不属于 \(\mu\)-本质像（\(g\) 的）。

反之，若 \(z \in \mathbf{C}\) 附着于 \(g(Y)\)，则对 z 的每个开邻域 \(U\) ，集 \(g^{-1}(U)\) 是 X 中与 Y 相交的开集。因此它不是 \(\mu\)-零测的，又由于它是开的，它不是局部 \(\mu\)-零测的（INT, IV, p. 172, § 5, n°2，推论 2）。于是 z 属于 g 的 \(\mu\)-本质像。

引理 3. --- 设 \(g\) 是从 X 到 \(\mathbf{C}\) 中的连续函数，S 是它的 \(\mu\)-本质像。假设 S 非空且 \(0 \notin S\) ，并用 \(\delta\) 表示 0 到 S 的距离。则 \(\delta > 0\)。

置 \(h(x) = 1 / g(x)\) （若 \(g(x) \neq 0\)），否则置 \(h(x) = 0\) 。函数 \(h\) 属于 \(\mathcal{L}^{\infty}(\mathrm{X},\mu)\)。设 \(\widetilde{h}\) 是 \(h\) 在 \(\mathrm{L}^{\infty}(\mathrm{X},\mu)\) 中的类。我们有公式 \(\delta^{-1} = \| \widetilde{h}\|_{\infty}\) 。

设 U 是 0 的开邻域，使得开集 \(Z = g^{-1}(U)\) 局部 \(\mu\)-零测，从而是零测的（INT, IV, p. 172, § 5, no. 2，推论 2）。设 Y 是 \(\mu\) 的支撑；我们有 \(Y \subset X - Z\) 。函数 \(h\) 在 \(X - Z\) 上的限制连续且有界，故 \(h \in \mathcal{L}^{\infty}(X, \mu)\) ，且 \(\widetilde{h}\) 在 \(\mathcal{L}^{\infty}(X, \mu)\) 中的范数等于它在 \(X - Z\) 上的限制的范数。此外，对每个 \(\alpha \in R_{+}\) ，\(x \in X - Z\) 中满足 \(|h(x)| > \alpha\) 者所成之集在 \(X - Z\) 中是开的；因此它局部 \(\mu\)-零测当且仅当它不与 Y 相交（INT, IV，同前及 INT, III, p. 66, § 3, no. 2，定义 1）。于是，我们有

\[
\| \widetilde {h} \| _ {\infty} = \sup _ {x \in \mathrm{Y}} \frac {1}{| g (x) |},
\]

由此得

\[
\frac {1}{\| \widetilde {h} \| _ {\infty}} = \inf _ {x \in \mathrm{Y}} | g (x) | = \inf _ {\lambda \in g (\mathrm{Y})} | \lambda | = \inf _ {\lambda \in \overline {{g (\mathrm{Y})}}} | \lambda |
\]

由前一引理，它等于 δ。

\subsection{普遍可测函数}

在本节中，X 是局部紧拓扑空间。我们回顾（INT, V, p. 28, § 3, n° 4，定义 2），从 X 到拓扑空间 Y 中的映射 f 称为普遍可测的，如果它对 X 上的每个正测度 \(\mu\) 都是 \(\mu\)-可测的。为此只需 f 对 X 上每个具有紧支撑的正测度 \(\mu\) 是 \(\mu\)-可测的（INT, V, p. 28, § 4, no. 3，命题 6）。

引理 4. --- 设 Y 与 Z 是局部紧拓扑空间，\(f\colon X\to Y\) 与 \(g\colon Y\to Z\) 是普遍可测映射。映射 \(g\circ f\colon X\to Z\) 是普遍可测的。

设 \(\mu\) 是 X 上具有紧支撑的正测度，C 是它的支撑。\(f\) 在 C 上的限制是 \((\mu|\mathrm{C})\)-正常的。映射 \(g\) 关于像测度 \((f|\mathrm{C})(\mu|\mathrm{C})\) 可测，故映射 \((g \circ f)|\mathrm{C} = g \circ (f|\mathrm{C})\) 是 \((\mu|\mathrm{C})\)-可测的。因此映射 \(g \circ f\) 是 \(\mu\)-可测的。引理得证。

我们用 \(\mathcal{L}_{\mathrm{u}}(\mathrm{X})\) 表示在 X 上普遍可测的复值函数所成的向量空间，用 \(\mathcal{L}_{\mathrm{u}}^{\infty}(\mathrm{X})\) 表示在 X 上有界的函数 \(f \in \mathcal{L}_{\mathrm{u}}(\mathrm{X})\) 所成的子空间。对 \(f \in \mathcal{L}_{\mathrm{u}}^{\infty}(\mathrm{X})\)，我们记 \(\|f\|_{\infty} = \sup_{x \in X} |f(x)|\) 。

命题 3. --- a) 集 \(\mathcal{L}_{\mathrm{u}}(\mathrm{X})\) 是从 X 到 C 中的函数所成的对合代数的对合子代数；

\begin{enumerate}
\def\labelenumi{\alph{enumi})}
\setcounter{enumi}{1}
\item
  集 \(\mathcal{L}_{\mathrm{u}}^{\infty}(\mathrm{X})\) 是 \(\mathcal{L}_{\mathrm{u}}(\mathrm{X})\) 的子代数，且由 \(f\mapsto \| f\|_{\infty}\) 定义的映射是 \(\mathcal{L}_{\mathrm{u}}^{\infty}(\mathrm{X})\) 上的一个范数，对此范数 \(\mathcal{L}_{\mathrm{u}}^{\infty}(\mathrm{X})\) 是有单位元的星代数；
\item
  代数 \(\mathcal{L}_{\mathrm{u}}(\mathrm{X})\) 包含 \(\mathcal{C}(\mathrm{X})\) ，且 \(\mathcal{L}_{\mathrm{u}}^{\infty}(\mathrm{X})\) 包含 \(\mathcal{C}_b(\mathrm{X})\) ；
\item
  从 X 到 \(\mathbf{C}\) 中的每个博雷尔函数（TG, IX, p. 61，定义 5）都属于 \(\mathcal{L}_{\mathrm{u}}(\mathrm{X})\) ；
\item
  对任意序列 \((f_n)_{n\in \mathbf{N}}\) ——\(\mathcal{L}_{\mathrm{u}}(\mathrm{X})\) 中逐点收敛到函数 \(f\) （从 X 到 \(\mathbf{C}\) 中的函数）的序列——我们有 \(f\in \mathcal{L}_{\mathrm{u}}(\mathrm{X})\) 。
\end{enumerate}

断言 \(a\)）与 \(c\)）由定义得出（参见 INT, IV, p. 175, § 5, no. 3，推论 3）。断言 \(d\)）是 INT, IV, p. 179, § 5, no. 5 定理 4 的推论，因为 \(\mathbf{C}\) 的每个博雷尔子集在博雷尔函数下的原像是 X 的博雷尔子集，从而是普遍可测的（INT, V, p. 28, § 3, no. 4）。最后，断言 \(e\)）由叶戈罗夫定理得出（INT, IV, p. 175, § 5, no. 4，定理 2）。

为证明断言 \(b\)），只需注意断言 \(e\)）蕴含，\(a\) 更不用说，普遍可测函数的一致极限是普遍可测的。

X 上可能存在不是博雷尔的普遍可测函数。事实上，若 X 可度量化，X 的苏斯林子集的特征函数是普遍可测的（参见 INT, IV, p. 171, § 5, no. 1，推论 2）；然而在每个不可数的波兰空间中，都存在不是博雷尔的苏斯林子集（“苏斯林定理”；这由 TG, IX, p. 120 习题 8 以及对每个不可数波兰空间 X 存在逆映射为博雷尔的双射 \(\mathrm{X}\to\mathbf{N}^{\mathbf{N}}\) 这一事实得出）。

引理 5. --- 设 \(\mu\) 是 X 上的正测度。设 \(g\) 是 \(\mu\)-可测函数（从 X 到 \(\mathbf{C}\) 中的），并设 \(S\) 是它的 \(\mu\)-本质像。对每个函数 \(f \in \mathcal{L}_{\mathrm{u}}(\mathrm{S})\)，函数 \(h\) （从 X 到 \(\mathbf{C}\) 中的函数），若满足 \(h(x) = 0\) （当 \(g(x) \notin \mathrm{S}\)）与 \(h(x) = f(g(x))\) （当 \(g(x) \in \mathrm{S}\)），则是 \(\mu\)-可测的。

设 \(x \in X\) ，\(U\) 是 x 的相对紧开邻域。集 \(\mathrm{N} = \mathrm{U} - (\mathrm{U} \cap g^{-1}(\mathrm{S}))\) 是 \((\mu|\mathrm{U})\)-零测的。对每个整数 \(n \geqslant 1\)，设 \(V_{n}\) 是满足 \(N \subset V_{n} \subset U\) 且 \(\mu(V_{n} - N) < 1/n\) 的开集（INT, IV, p. 116, § 1, no. 4，命题 19）。

限制 \(\widetilde{g}\) （\(g\) 在 \(\mathrm{U}-\mathrm{V}_n\) 上的限制）是 \(\mu|(\mathrm{U}-\mathrm{V}_n)\)-正常的，且其像含于 S。由于 \(f\) 普遍可测，映射 \(x\mapsto f(g(x))\) （从 \(\mathrm{U}-\mathrm{V}_n\) 到 \(\mathbf{C}\) 中的映射）关于测度 \(\mu|(\mathrm{U}-\mathrm{V}_n)\) 可测（INT, V, p. 71, § 6, no. 2，命题 3）。因此，映射 \(h_n\) （从 X 到 \(\mathbf{C}\) 中的映射）——若 \(h_n(x) = 0\) （当 \(x\notin \mathrm{U}-\mathrm{V}_n\)）且 \(h_n(x) = f(g(x))\) （当 \(x\in \mathrm{U}-\mathrm{V}_n\)）——是 \(\mu\)-可测的（INT, IV, p. 193, § 5, no. 10，命题 16）。

由于 \(h_{n}(x) \to h(x)\) 对 \(\mu\)-几乎所有的 \(x \in U\) 成立，\(h\) 在 U 上的限制是 \(\mu\)-可测的（INT, IV, p. 175, § 5, no. 4，定理 2）。于是可断定 \(h\) 是 \(\mu\)-可测的，依据是局部化原理（INT, IV, p. 171, § 4, no. 2，命题 4）。

注. --- 在该引理的假设下，我们约定滥用记号，用 \(f \circ g\) 表示引理中定义的函数 \(h\) 。

若 \(g = g_{1}\) 在 X 上局部 \(\mu\)-几乎处处成立，则函数 \(g\) 与 \(g_{1}\) 有相同的 \(\mu\)-本质像，且有 \(f \circ g = f \circ g_{1}\) 。

若 \(g\) 是在 X 上 \(\mu\)-几乎处处定义的 \(\mu\)-可测函数，我们也用 \(f \circ g\) 表示函数 \(f \circ g_{1}\) ，其中 \(g_{1}\) 是定义在 X 上且局部 \(\mu\)-几乎处处等于 g 的 \(\mu\)-可测函数。

\subsection{\texorpdfstring{星代数 \(\mathrm{L}^{\infty}(\mathrm{X},\mu)\)}{星代数 \textbackslash mathrm\{L\}\^{}\{\textbackslash infty\}(\textbackslash mathrm\{X\},\textbackslash mu)}}

设 X 是局部紧拓扑空间，\(\mu\) 是 X 上的正测度。我们考虑交换星代数 \(\mathrm{L}^{\infty}(\mathrm{X},\mu)\) （I, p. 103 的例 4）。

引理 6. --- 设 \(g \in \mathcal{L}^{\infty}(\mathrm{X}, \mu)\)。\(g\) 在 \(\mathrm{L}^{\infty}(\mathrm{X}, \mu)\) 中的类的谱等于 \(\mu\)-本质像（\(g\) 的）。

用 \(\widetilde{g}\) 表示 \(g\) 在 \(\mathrm{L}^{\infty}(\mathrm{X},\mu)\) 中的类，用 \(S\) 表示 \(\mu\)-本质像（\(g\) 的）。设 \(z\in \mathbf{C} - \mathbf{S}\)。按定义，存在开邻域 \(U\) （\(z\) 的），使得 \(\mathrm{Y} = g^{-1}(\mathrm{U})\) 局部 \(\mu\)-零测。函数 \(h\) ——由 \(h(x) = (g(x) - z)^{-1}\) （当 \(x\notin \mathrm{Y}\)）与 \(h(x) = 0\) （当 \(x\in \mathrm{Y}\)）定义——属于 \(\mathcal{L}^{\infty}(\mathrm{X},\mu)\)。它的类 \(\widetilde{h}\) 在 \(\mathrm{L}^{\infty}(\mathrm{X},\mu)\) 中满足 \((\widetilde{g} -z)\widetilde{h} = 1\) ，因为 \((g(x) - z)h(x) = 1\) 对每个 \(x\) （只要它不属于局部 \(\mu\)-零测集 Y）成立。因此 \(z\in \mathbf{C} - \mathrm{Sp}(\widetilde{g})\) 。

反之，设 \(z \in \mathbf{C} - \mathrm{Sp}(\widetilde{g})\)。设 \(h \in \mathcal{L}^{\infty}(\mathrm{X},\mu)\) 是这样的函数：它的类是 \(\widetilde{g} - z\) 在 \(\mathrm{L}^{\infty}(\mathrm{X},\mu)\) 中的逆。存在实数 \(\mathrm{M} > 0\) ，使得 \(|h(x)| \leqslant \mathrm{M}\) 局部 \(\mu\)-几乎处处成立，此外我们有 \((g(x) - z)h(x) = 1\) 局部 \(\mu\)-几乎处处成立。设 U 是以 \(z\) 为心、\(\mathrm{M}^{-1}\) 为半径的开球（在 \(\mathbf{C}\) 中）；则 \(g^{-1}(\mathrm{U})\) 含于局部 \(\mu\)-零测集

\[
h ^ { - 1 } ( ] \mathrm{M} , + \infty [ ) \cup \{ x \in \mathrm{X} \mid ( g ( x ) - z ) h ( x ) \neq 1 \} ,
\]

从而 \(z \in \mathbf{C} - \mathrm{S}\) 。引理得证。

命题 4. --- 设 \(g \in \mathcal{L}^{\infty}(\mathrm{X},\mu)\) 。用 \(\widetilde{g}\) 表示 \(g\) 在 \(\mathrm{L}^{\infty}(\mathrm{X},\mu)\) 中的类，用 \(\mathrm{S}\) 表示 \(\widetilde{g}\) 的谱。映射 \(f \mapsto f \circ g\) （p. 184 的注）是从 \(\mathcal{C}(\mathrm{S})\) 到 \(\mathrm{L}^{\infty}(\mathrm{X},\mu)\) 中的与 \(\widetilde{g}\) 相伴的泛函演算映射。

设 \(f \in \mathcal{C}(\mathrm{S})\) 。函数 \(h = f \circ g\) 是 \(\mu\)-可测的（IV, p. 184 的引理 5）且有界。用 \(\widetilde{f \circ g}\) 表示它在 \(\mathrm{L}^{\infty}(\mathrm{X},\mu)\) 中的类。映射 \(f \mapsto \widetilde{f \circ g}\) 是从 \(\mathcal{C}(\mathrm{S})\) 到 \(\mathrm{L}^{\infty}(\mathrm{X},\mu)\) 中的连续有单位元的对合代数同态。由于 \(g(x)\) 属于 \(\mathrm{S}\) （局部 \(\mu\)-几乎处处；IV, p. 181 的引理 1 及引理 6），这个同态把 \(\mathrm{S}\) 上的恒等映射映到类 \(\widetilde{g}\) （函数 \(g\) 的类）。因此它与 \(\widetilde{g}\) 的泛函演算映射相合（I, p. 111 的命题 7）。